\chapter{Bedford--Taylor products and local capacity}
\label{chap:q08}

On a domain in $\mathbb{C}^n$, two currents with measure coefficients cannot in general be multiplied. A locally bounded plurisubharmonic potential can be multiplied by a measure, however. Applying $\ddc$ to that product gives a positive closed current. Iterating this construction produces the local Monge--Amp\`ere operator. Its continuity along decreasing potentials supplies the symmetry which is absent from the recursive definition, and its measures lead to a capacity adapted to complex poles.

\section{Bounded potentials and recursive products}
\label{sec:bt-construction}

Let $\Omega\subset\mathbb{C}^n$, $n\ge1$, be open, and let $T$ be a positive closed current of bidimension $(r,r)$, $0\le r\le n$. Thus its bidegree is $(n-r,n-r)$. Put $\beta=(\mathrm{i}/2)\sum_j\mathrm{d}z_j\wedge\mathrm{d}\bar z_j$, as in \cref{chap:complex-currents}, and write
\[
 \|T\|_K=\int_KT\wedge\frac{\beta^r}{r!}
\]
for its trace mass on a compact set $K$. The coefficient estimates of \cref{thm:positive-current-order-zero} compare this mass with all coefficient variations. Local boundedness below as well as above is required for the potentials in this section.

\begin{definition}\label{def:bedford-taylor-product}
For a locally bounded plurisubharmonic $u$ and $r\ge1$, define
\[
 \ddc u\wedge T=\ddc(uT).
\]
Here $uT$ means multiplication of the coefficient measures by the bounded Borel representative $u$. For $q\le r$, define recursively
\[
 \ddc u_1\wedge\cdots\wedge\ddc u_q\wedge T
 =\ddc\bigl(u_1\ddc u_2\wedge\cdots\wedge\ddc u_q\wedge T\bigr).
\]
The empty product is $T$. With $T=1$ and $q=n$, the resulting top-degree current is the Monge--Amp\`ere measure $(\ddc u)^n$.
\end{definition}
\index{Bedford--Taylor product}\index{Monge--Amp\`ere measure!locally bounded potentials}

\begin{proposition}\label{prop:bt-positive-recursion}
Every product in \cref{def:bedford-taylor-product} is positive and closed. For smooth potentials it agrees with the exterior product. The construction is local, is additive in each potential and in the base current whenever the summands satisfy the hypotheses, and is invariant under biholomorphic coordinates.
\end{proposition}
\begin{proof}
For one potential, choose local radial regularizations $u_\varepsilon\downarrow u$ from \cref{thm:psh-distribution-representative}. They are uniformly bounded on each smaller compact set. Dominated convergence against the coefficient variations of $T$ gives $u_\varepsilon T\to uT$ weakly. Since $T$ is closed and of pure bidegree, $\partial T=\bar\partial T=0$; the distributional product rule for a smooth multiplier therefore gives
\[
 \ddc(u_\varepsilon T)=\ddc u_\varepsilon\wedge T.
\]
The right side is positive and closed. Differentiation is weakly continuous, so its limit is the current in the definition and has the same properties. Recursion proves the assertion for all $q$.

For smooth potentials the same product rule proves agreement at every step. Locality and additivity follow directly from multiplication and differentiation; no product of two measure coefficients is taken. A biholomorphism pulls coefficient measures back and preserves their multiplication by a Borel function. Its smooth pullback commutes with $\partial,\bar\partial$, and hence with their distributional extensions by coordinate changes on tests. The recursive formulas consequently agree in overlapping holomorphic charts.
\end{proof}

For $u\in C^2$, determinant expansion and $\mathrm{i}\,\mathrm{d}z_j\wedge\mathrm{d}\bar z_j=2\,\mathrm{d}x_j\wedge\mathrm{d}y_j$ give
\begin{equation}\label{eq:bt-determinant-normalization}
 (\ddc u)^n=\frac{n!}{\pi^n}\det(u_{j\bar k})\,\mathrm{d}V_{2n}.
\end{equation}
In particular $\ddc|z|^2=\beta/\pi$. In dimension zero the empty product is $1$, and integration uses counting measure; there is no positive-degree product.

\begin{theorem}[Chern--Levine--Nirenberg]\label{thm:cln-bounded}
For $K\Subset L^\circ\subset L\Subset\Omega$ and $q\le r$, there is a constant depending only on these sets, the dimension and $q$ such that
\[
 \|\ddc u_1\wedge\cdots\wedge\ddc u_q\wedge T\|_K
 \le C_{K,L,q}\|T\|_L\prod_{j=1}^q\|u_j\|_{L^\infty(L)}.
\]
The bounded norms refer to the specified pointwise representatives.
\end{theorem}
\begin{proof}
For $q=1$, take a nonnegative smooth cutoff $\chi$ supported in $L^\circ$ and equal to one near $K$. Positivity, followed by distributional integration by parts, gives
\[
 \begin{aligned}
 \int_K\ddc u\wedge T\wedge\beta^{r-1}
 &\le\int\chi\,\ddc(uT)\wedge\beta^{r-1}\\
 &=\int uT\wedge\ddc\chi\wedge\beta^{r-1}
 \le C\|u\|_{L^\infty(L)}\|T\|_L.
 \end{aligned}
\]
The last bound uses the coefficient variations controlled by the trace; the signed form $\ddc\chi$ need not be positive. Such a cutoff is supplied by \cref{prop:continuous-density}. For $q>1$, choose $q-1$ intermediate compact neighborhoods between $K$ and $L^\circ$ and apply this estimate successively to the positive closed recursive products. The factorials in the trace norms are absorbed into the stated constant.
\end{proof}
\index{Chern--Levine--Nirenberg inequality}

\begin{proposition}\label{prop:bt-uniform-continuity}
Suppose $u_{j,k}\to u_j$ locally uniformly, where the limits are finite continuous plurisubharmonic functions and the approximating functions are locally bounded and plurisubharmonic. If $T_k\to T$ weakly for positive closed currents of fixed bidimension, then the corresponding products, and products with one additional uniformly convergent continuous-limit potential as weight, converge weakly.
\end{proposition}
\begin{proof}
Induct on the number of Hessian factors. Their local masses are bounded by \cref{thm:cln-bounded} and the trace bounds for $T_k$. To pass from $S_k\to S$ to $v_kS_k\to vS$, approximate the continuous $v$ uniformly on the test support by a smooth $v_\varepsilon$ and write
\[
 v_kS_k-vS=(v_k-v)S_k+(v-v_\varepsilon)S_k
             +v_\varepsilon(S_k-S)+(v_\varepsilon-v)S.
\]
The mass bounds control the first, second and fourth terms. For fixed $\varepsilon$ the third tends weakly to zero. First let $k$ increase and then $\varepsilon$ decrease. Applying $\ddc$ proves the next unweighted assertion. Smooth approximation of a continuous function follows from \cref{prop:continuous-density} in coordinate neighborhoods.
\end{proof}

\section{Decreasing continuity and symmetry}
\label{sec:bt-monotone}

Uniform convergence is unavailable for bounded discontinuous potentials. A common smooth collar permits an integration-by-parts comparison of total weighted mass. Together with an upper bound from upper semicontinuity, this identifies every weak limit.

For the calculations below put $\mathrm{d}^{\mathrm c}=(\partial-\bar\partial)/(4\pi\mathrm{i})$, so that $\mathrm{d}\mathrm{d}^{\mathrm c}=\ddc$. On real functions $a,b$ the mixed terms satisfy $\mathrm{d}a\wedge\mathrm{d}^{\mathrm c}b\wedge R=\mathrm{d}b\wedge\mathrm{d}^{\mathrm c}a\wedge R$ whenever the closed pure-type current $R$ makes the product top degree. This follows by expanding the two types; only the $(1,1)$ terms survive.

\begin{lemma}\label{lem:bt-common-collar}
Let $B$ be a ball whose closure lies in the domain of a positive closed current $R$ of bidimension $(1,1)$. Let $a$ be smooth and $b$ bounded plurisubharmonic near $\overline B$. If both equal the same smooth function $\rho$ on a collar of $\partial B$, then
\[
 \int_B a\,\ddc b\wedge R=\int_B b\,\ddc a\wedge R.
\]
The assertion also applies after wedging $R$ with any smooth closed complementary form that preserves positivity.
\end{lemma}
\begin{proof}
Write $a=\rho+\alpha$ and $b=\rho+h$, with $\alpha$ smooth and $h$ bounded, both compactly supported inside $B$. The current $\ddc(hR)=\ddc b\wedge R-\ddc\rho\wedge R$ has measure coefficients. Distributional integration by parts, tested with cutoffs equal to one near the compact supports, gives
\[
 \int\rho\,\ddc(hR)=\int h\,\ddc\rho\wedge R,
 \qquad
 \int\alpha\,\ddc(hR)=\int h\,\ddc\alpha\wedge R.
\]
The same identity for the smooth $\alpha$ interchanges it with $\rho$. Expanding the two integrals in the statement now cancels their common $\rho\,\ddc\rho\wedge R$ term and pairs every remaining term. All integrals are finite by boundedness and the local trace estimates. This proof uses compact supports of the differences, rather than a boundary value assertion for nonsmooth derivatives.
\end{proof}

\begin{lemma}\label{lem:usc-weight-weak-limit}
Let $f_j$ be locally bounded real upper semicontinuous functions decreasing to a locally bounded $f$. If $\mu_j\to\mu$ weakly for positive measures and $f_j\mu_j\to\nu$ along a subsequence, then $\nu\le f\mu$ as measures.
\end{lemma}
\begin{proof}
Fix a nonnegative compactly supported test. For a fixed index $l$, $f_j\le f_l$ eventually. On a compact neighborhood of the support choose continuous majorants $g_m\downarrow f_l$ from \cref{lem:usc-majorants}. Weak convergence gives $\int\chi\,\mathrm{d}\nu\le\int\chi g_m\,\mathrm{d}\mu$. Bounded convergence as $m\to\infty$ gives the bound with $f_l$; bounded convergence again as $l\to\infty$ gives the asserted bound with $f$. Local boundedness and weak convergence bound all measures on the chosen compact neighborhood, justifying these tests.
\end{proof}

\begin{lemma}\label{lem:smooth-maximum}
There are smooth convex functions $M_\varepsilon:\mathbb{R}^2\to\mathbb{R}$, increasing in each coordinate, such that $M_\varepsilon\downarrow\max$ as $\varepsilon\downarrow0$, $\max\le M_\varepsilon\le\max+\varepsilon$, and $M_\varepsilon(s,t)=\max\{s,t\}$ when $|s-t|>2\varepsilon$. If $a,b$ are smooth plurisubharmonic functions, so is $M_\varepsilon(a,b)$.
\end{lemma}
\begin{proof}
Convolve $\max\{s,t\}$ with a nonnegative smooth even radial probability supported in the radius-$\varepsilon$ ball of $\mathbb{R}^2$. Convexity and coordinate monotonicity are preserved. Jensen's inequality gives the lower bound; the Lipschitz bound for the maximum gives the upper bound. Where $|s-t|>2\varepsilon$, the same linear coordinate is maximal throughout that ball, and its mean equals its center value. For fixed $(s,t)$, the average as a function of $\varepsilon$ is convex and even, hence increasing for $\varepsilon\ge0$. This proves the decreasing assertion. Finally the Hessian chain rule has a nonnegative sum of the Hessians of $a,b$, plus the nonnegative real Hessian of $M_\varepsilon$ applied to their first derivatives, proving plurisubharmonicity.
\end{proof}

\begin{theorem}[Bedford--Taylor]\label{thm:bt-decreasing-continuity}
Let $u_{l,j}\downarrow u_l$, $0\le l\le q\le r$, be decreasing sequences of locally bounded plurisubharmonic functions with locally bounded limits. For a fixed positive closed $T$ of bidimension $(r,r)$,
\[
 \begin{aligned}
 \ddc u_{1,j}\wedge\cdots\wedge\ddc u_{q,j}\wedge T
 &\longrightarrow\ddc u_1\wedge\cdots\wedge\ddc u_q\wedge T,\\
 u_{0,j}\ddc u_{1,j}\wedge\cdots\wedge\ddc u_{q,j}\wedge T
 &\longrightarrow u_0\ddc u_1\wedge\cdots\wedge\ddc u_q\wedge T
 \end{aligned}
\]
weakly. The product is symmetric in its Hessian factors.
\end{theorem}
\begin{proof}
We prove the weighted assertion by induction on $q$, allowing every fixed positive closed base current of the required bidimension. For $q=0$ it is dominated convergence against the coefficient variations of $T$. Suppose the weighted assertion holds with $q-1$ Hessian factors. Applying $\ddc$ to that assertion proves the unweighted assertion with $q$ factors. Applying it to simultaneous smooth radial regularizations proves symmetry for $q$ factors, since the smooth exterior product is symmetric. These two conclusions are therefore available before proving the weighted assertion with $q$ factors.

The proof is local. Fix a ball $B=B(a,R)$ with closure in a slightly larger coordinate ball and a smaller compact set where the assertion is to be tested. Subtract constants so that all the potentials in question lie between $-M$ and $-1$ on the larger ball. Put $\psi=|z-a|^2-R^2$. Choose $\delta>0$ so that the smaller compact set lies in $\{\psi<-\delta\}$ and set $A=M/\delta$. Replacing every potential $v$ by $\max\{v,A\psi\}$ leaves it unchanged on that compact set, preserves the decreasing sequences, and makes all functions equal to $\rho=A\psi$ on one fixed collar of $\partial B$. Thus it suffices to treat this common-collar situation.

Every fixed modified $u_{l,j}$ has decreasing smooth plurisubharmonic approximations with exactly the same collar $\rho$: take
$M_\varepsilon(u_{l,j}*\rho_\varepsilon,A\psi)$, where the first convolution is performed on the original unmodified potential on the larger ball. For sufficiently small $\varepsilon$ these functions equal $A\psi$ on a common collar, decrease to the modified potential, and remain negative on $B$. Indeed the original smoothing is at most $-1$, the transition can occur only where $A\psi$ is bounded away from zero, and \cref{lem:smooth-maximum} is exact on the remaining collar. Denote these approximations by $u_{l,j}^{\varepsilon}$.

Write $S_j=\ddc u_{1,j}\wedge\cdots\wedge\ddc u_{q,j}\wedge T$ and $S=\ddc u_1\wedge\cdots\wedge\ddc u_q\wedge T$. Their convergence is already proved. The CLN estimate bounds the masses of $u_{0,j}S_j$ on compact sets. The coefficient-measure selection theorem \cref{thm:local-measure-selection}, applied to real and imaginary variations or to positive and negative parts, supplies weakly convergent subsequences. Let $\Theta$ be any such limit. Wedging with a smooth strongly positive test form and applying \cref{lem:usc-weight-weak-limit} gives $\Theta\le u_0S$, meaning that $u_0S-\Theta$ is a positive current. All these currents agree on the fixed collar, so the difference has compact support inside $B$.

Set $\gamma=\beta^{r-q}$ and
\[
 I=\int_Bu_0S\wedge\gamma,
 \qquad I_j=\int_Bu_{0,j}S_j\wedge\gamma.
\]
For each fixed $j$ and each collection of small smoothing parameters, successively replace the undifferentiated $u_l$ by its larger $u_{l,j}^{\varepsilon_l}$ and interchange it with the next undifferentiated potential using \cref{lem:bt-common-collar}. Every differentiated current in this process is positive and closed. This gives
\[
 I\le\int_Bu_{q,j}^{\varepsilon_q}
       \ddc u_{0,j}^{\varepsilon_0}\wedge\cdots\wedge
       \ddc u_{q-1,j}^{\varepsilon_{q-1}}\wedge T\wedge\gamma.
\]
For example the first two steps are
$I\le\int u_{0,j}^{\varepsilon_0}\ddc u_1\wedge\cdots
=\int u_1\ddc u_{0,j}^{\varepsilon_0}\wedge\ddc u_2\wedge\cdots$;
positivity then permits the next replacement. This explains the direction of each inequality even though the weights are negative.

Let $\varepsilon_q$ decrease to zero first, then $\varepsilon_{q-1},\ldots,\varepsilon_1$ in that order, keeping $\varepsilon_0$ fixed. At each step there are at most $q-1$ nonsmooth Hessian factors with a bounded weight. The induction hypothesis, with the remaining smooth factors absorbed into the fixed base current, gives convergence of these integrals. A cutoff can replace integration over $B$, since all currents in question have a fixed collar. The resulting integral is
\[
 \int_Bu_{q,j}\ddc u_{0,j}^{\varepsilon_0}\wedge
       \ddc u_{1,j}\wedge\cdots\wedge\ddc u_{q-1,j}\wedge T\wedge\gamma.
\]
The common-collar identity interchanges the smooth $u_{0,j}^{\varepsilon_0}$ and the bounded $u_{q,j}$; symmetry of the $q$-factor product, already established above, then puts its Hessians in their original order. The expression becomes $\int_Bu_{0,j}^{\varepsilon_0}S_j\wedge\gamma$. Dominated convergence against the fixed $S_j\wedge\gamma$ now lets $\varepsilon_0\downarrow0$ and proves $I\le I_j$.

Along the subsequence defining $\Theta$, the common collar makes $I_j\to\int_B\Theta\wedge\gamma$. Thus $\int_B(u_0S-\Theta)\wedge\gamma\le0$. Positivity and the trace coefficient bound in \cref{thm:positive-current-order-zero} force $u_0S=\Theta$. Every weak subsequential limit is therefore the desired one. Local relative compactness gives convergence of the full sequence. This completes the induction and the theorem.
\end{proof}
\index{Bedford--Taylor product!decreasing continuity}\index{Bedford--Taylor product!symmetry}

\begin{corollary}\label{cor:bt-polynomial-weights}
In \cref{thm:bt-decreasing-continuity}, one may replace the single weight by any real polynomial in finitely many bounded plurisubharmonic potentials, each approximated by a decreasing sequence with its bounded limit. The common-collar identity of \cref{lem:bt-common-collar} also holds when both $a$ and $b$ are bounded plurisubharmonic.
\end{corollary}
\begin{proof}
Shift the potentials by constants to make them nonnegative on a compact neighborhood of the test support. Every monomial is a linear combination of powers of their nonnegative sums, by
\[
 \prod_{l=1}^N a_l=\frac{1}{N!}
   \sum_{I\subset\{1,\ldots,N\}}(-1)^{N-|I|}
        \left(\sum_{l\in I}a_l\right)^N.
\]
Expansion proves this identity: a monomial missing a variable cancels in the alternating sum, and the coefficient of the product containing all $N$ variables is $N!$. Repeated variables are allowed. A nonnegative sum raised to the integer power $N$ is plurisubharmonic, by the convex increasing chain rule, and its smooth approximations decrease. The weighted continuity theorem therefore applies to every term and hence to the polynomial.

For the collar identity, radially smooth $a$ and $b$ with the same radius. Near the boundary they coincide with the same smooth $\rho*\rho_\varepsilon$. Apply the smooth identity and let the radius decrease, using weighted continuity with one Hessian factor. On a smaller fixed collar the approximated currents converge in coefficient variation, since their smooth multipliers and derivatives converge uniformly there. Thus a compactly supported test gives convergence of the interior integrals, while this variation convergence gives convergence on the remaining collar, including when the trace of $R$ charges the sphere. The integrals over $B$ consequently converge, proving the bounded identity.
\end{proof}

\begin{proposition}\label{prop:bt-maximum-locality}
For locally bounded plurisubharmonic $u,v$ and $p\le r$,
\[
 \mathbf{1}_{\{u>v\}}(\ddc\max\{u,v\})^p\wedge T
 =\mathbf{1}_{\{u>v\}}(\ddc u)^p\wedge T.
\]
The indicators restrict the coefficient measures of these order-zero currents. The set $\{u>v\}$ is Borel; Euclidean openness is not assumed.
\end{proposition}
\begin{proof}
Work locally and choose smooth $u_j\downarrow u$. Since $v$ is upper semicontinuous, $\{u_j>v\}$ is Euclidean open. On it $w_j=\max\{u_j,v\}$ equals $u_j$, so the recursive construction is local there. The nonnegative Borel weight $f_j=w_j-v$ vanishes on its complement, giving
\[
 f_j(\ddc w_j)^p\wedge T=f_j(\ddc u_j)^p\wedge T.
\]
The functions $w_j$ decrease to $w=\max\{u,v\}$. Apply weighted decreasing continuity separately with the weights $w_j$ and $v$ on each side. The limit is
$f(\ddc w)^p\wedge T=f(\ddc u)^p\wedge T$, where $f=w-v\ge0$.
Multiply its coefficient measures by $1/(f+\varepsilon)$, then let $\varepsilon\downarrow0$. Dominated convergence for their locally finite variations gives the indicator of $\{f>0\}=\{u>v\}$, as asserted.
\end{proof}
\index{Bedford--Taylor product!maximum locality}

\begin{theorem}\label{thm:bt-comparison}
Let $D\subset\mathbb{C}^n$ be bounded and let $u,v$ be bounded plurisubharmonic functions on $D$. If $\liminf_{D\ni z\to x}(u(z)-v(z))\ge0$ for every $x\in\partial D$, then
\[
 \int_{\{u<v\}}(\ddc v)^n\le\int_{\{u<v\}}(\ddc u)^n.
\]
The integrals may be infinite. If $(\ddc u)^n=0$, the boundary condition implies $u\ge v$ everywhere.
\end{theorem}
\begin{proof}
For $\varepsilon>0$, compactness of $\partial D$ and the boundary lower limits show that $w=\max\{u,v-\varepsilon\}$ equals $u$ outside a compact subset of $D$. Symmetry and additivity of the bounded products give
\[
 (\ddc w)^n-(\ddc u)^n
 =\ddc\left((w-u)\sum_{l=0}^{n-1}
           (\ddc w)^l\wedge(\ddc u)^{n-1-l}\right).
\]
The current inside the parentheses has compact support. Test with a cutoff equal to one on that support to conclude that the difference of the measures has total mass zero. Choose a relatively compact open $D'$ containing that support and $\{u\le v-\varepsilon\}$; both measures are finite on $D'$. Maximum locality gives
\[
 \begin{aligned}
 \int_{\{u<v-\varepsilon\}}(\ddc v)^n
 &\le\int_{D'}(\ddc w)^n
       -\int_{D'\cap\{u>v-\varepsilon\}}(\ddc u)^n\\
 &=\int_{\{u\le v-\varepsilon\}}(\ddc u)^n
 \le\int_{\{u<v\}}(\ddc u)^n.
 \end{aligned}
\]
The sets on the left increase to $\{u<v\}$ as $\varepsilon\downarrow0$. Monotone convergence proves the comparison.

For the last assertion choose $R$ with $D\subset B(0,R)$ and replace $v$ by $v_\eta=v+\eta(|z|^2-R^2)$. The boundary condition still holds and $(\ddc v_\eta)^n\ge\eta^n(\beta/\pi)^n$. Comparison makes $\{u<v_\eta\}$ Lebesgue null. An almost-everywhere inequality between two plurisubharmonic functions holds everywhere: their maximum agrees almost everywhere with the larger canonical representative, so uniqueness in \cref{thm:psh-distribution-representative} identifies it pointwise. Thus $u\ge v_\eta$ everywhere; let $\eta\downarrow0$.
\end{proof}
\index{comparison principle!bounded plurisubharmonic functions}

\begin{corollary}\label{cor:bt-local-domination}
Let $D\subset\mathbb C^n$ be a bounded domain, $n\geq1$, and let $u,v$ be locally bounded plurisubharmonic functions. Suppose that
\[
 \liminf_{D\ni z\to x}(v(z)-u(z))\geq0
       \quad(x\in\partial D),\qquad
 u\leq v\quad\text{almost everywhere for }(\ddc v)^n.
\]
Then $u\leq v$ everywhere on $D$.
\end{corollary}
\begin{proof}
Choose $R$ with $D\subset B(0,R)$, and for $\varepsilon,\delta>0$ put
$u_{\varepsilon,\delta}=u-\varepsilon+\delta(|z|^2-R^2)$.
The boundary assumption and a finite cover of the compact boundary show that
$v>u-\varepsilon/2$ on a neighborhood of $\partial D$ inside $D$. Hence
$A=\{v<u_{\varepsilon,\delta}\}$ lies in a compact subset of $D$.
Choose a relatively compact open $D'\subset D$ containing that compact subset, with its boundary in the same neighborhood. Both potentials are bounded on $D'$, and $v>u_{\varepsilon,\delta}$ near its boundary. Comparison gives
\[
 \delta^n\int_A(\ddc|z|^2)^n
 \leq\int_A(\ddc u_{\varepsilon,\delta})^n
 \leq\int_A(\ddc v)^n=0.
\]
The last equality follows from $A\subset\{v<u\}$. Thus $u_{\varepsilon,\delta}\leq v$ almost everywhere, then everywhere by uniqueness of canonical psh representatives applied to their maximum. Let $\delta\downarrow0$ and $\varepsilon\downarrow0$.
\end{proof}
\index{domination principle!locally bounded potentials}

\section{Weighted estimates and Monge--Amp\`ere capacity}
\label{sec:bt-capacity}

A capacity takes the largest mass assigned to a set by a normalized class of Monge--Amp\`ere measures. Uniform integral estimates for that entire class are therefore needed before exceptional sets can be controlled.

\begin{theorem}\label{thm:cln-weighted}
For $K\Subset L^\circ\subset L\Subset\Omega$, $V\in\PSH(\Omega)$, and locally bounded plurisubharmonic $u_1,\ldots,u_q$, $0\le q\le n$,
\[
 \int_K |V|\,\ddc u_1\wedge\cdots\wedge\ddc u_q\wedge\beta^{n-q}
 \le C_{K,L,q}\|V\|_{L^1(L)}
            \prod_{l=1}^q\|u_l\|_{L^\infty(L)}.
\]
In particular the specified pointwise $V$ is integrable against every bounded-potential product trace. The base current here is $1$; an arbitrary singular $T$ is not covered by this $L^1$ estimate.
\end{theorem}
\begin{proof}
For $q=0$ this is the Lebesgue integral, up to a dimension constant. Otherwise cover $K$ by finitely many smaller balls inside balls with closure in $L^\circ$. If some bounded norm is zero, that potential is zero on the larger neighborhood and the local product vanishes. Divide the other potentials by twice their norms and subtract $3/2$, so they lie in $[-2,-1]$. The Hessians acquire only the product of those scaling constants. It suffices to prove a uniform estimate for the normalized functions on one smaller ball.

Use the common-collar modification from \cref{thm:bt-decreasing-continuity} on $B$, leaving the potentials unchanged near the smaller ball. Denote their product by $S$, put $\psi=|z-a|^2-R^2$ and $b=\ddc\psi=\beta/\pi$. There are $A,\tau>0$, depending only on the chosen balls, such that $S=A^q b^q$ on $B\setminus B_\tau$, where $B_\tau=\{\psi<-\tau\}$ contains the smaller ball. Choose a nonnegative smooth cutoff $\chi$ supported in $B$ and equal to $-\psi$ on a neighborhood of $\overline{B_\tau}$.

First take smooth $V$ with $V\le0$ on $B_\tau$, and suppose $q\le n-1$. Splitting the integral into $B_\tau$ and its collar gives the exact identity
\[
 \begin{aligned}
 \int_{B_\tau}(-V)S\wedge b^{n-q}
 &=\int_B V S\wedge b^{n-q-1}\wedge\ddc\chi
       -A^q\int_{B\setminus B_\tau}V b^{n-1}\wedge\ddc\chi\\
 &=\int_B\chi\,\ddc V\wedge S\wedge b^{n-q-1}
       -A^q\int_{B\setminus B_\tau}V b^{n-1}\wedge\ddc\chi.
 \end{aligned}
\]
For the first term, the bounded CLN estimate with base current $\ddc V$ gives a bound by its trace mass on a fixed larger compact set. Testing $\ddc V$ with a cutoff bounds that mass by $C\|V\|_{L^1(L)}$, exactly as in the one-factor CLN proof. The collar term is bounded by the same norm because it has smooth bounded coefficients. This proves the estimate for nonpositive smooth $V$.

For a general $V$, its ball submean inequality bounds its positive part on a fixed intermediate compact neighborhood by $C\|V\|_{L^1(L)}$. Its radial smoothings have the same type of bound and uniformly bounded $L^1$ norms on that neighborhood, by Fubini. Subtract this common constant and apply the smooth estimate to $V_\varepsilon$. The nonnegative weights $C\|V\|_{L^1(L)}-V_\varepsilon$ increase to the pointwise weight for $V$, so monotone convergence against the fixed positive product measure gives the estimate for its negative part. The bounded CLN estimate controls its positive part. This proves the result for $q\le n-1$.

For $q=n$, lift the functions to a product with a plane disc, independently of the extra variable, and apply the established result in dimension $n+1$. The trace of the $n$-factor product there includes exactly the area form of the extra variable, since its other terms repeat a top-degree $z$-factor and vanish. Fubini cancels its positive disc-area constant and gives the desired estimate on the original ball. The lifted $L^1$ norm has the same area factor. Combining the finite ball cover and reversing the normalizations proves the theorem.
\end{proof}
\index{Chern--Levine--Nirenberg inequality!weighted}

\begin{definition}\label{def:monge-ampere-capacity}
For a bounded open $D\subset\mathbb{C}^n$ and a Borel set $E\subset D$, put
\[
 \operatorname{Cap}_D(E)=\sup\left\{\int_E(\ddc v)^n:
            v\in\PSH(D),\ 0\le v\le1\right\}.
\]
For any subset $A\subset D$, its outer capacity is
\[
 \operatorname{Cap}_D^*(A)=\inf\{\operatorname{Cap}_D(G):A\subset G\subset D,
                                                  \ G\text{ open}\}.
\]
In local assertions we use relatively compact bounded coordinate domains, in particular balls.
\end{definition}
\index{capacity!Monge--Amp\`ere}\index{capacity!outer Monge--Amp\`ere}

\begin{proposition}\label{prop:bt-capacity-basic}
These capacities are increasing and countably subadditive. Capacity is finite on compact subsets of $D$, and for Borel $E$
$\operatorname{Cap}_D(E)=\sup_{K\subset E,\ K\text{ compact}}\operatorname{Cap}_D(K)$.
It is continuous on increasing Borel sequences. On each compact neighborhood Lebesgue measure is bounded by a constant times capacity. For a ball $B\Subset D$ and $K\Subset B$, there is a constant such that, for every Borel $E\subset K$,
\[
 \operatorname{Cap}_D(E)\le\operatorname{Cap}_B(E)
                              \le C\operatorname{Cap}_D(E).
\]
The corresponding local comparison also holds for outer capacities of subsets of $K$. Both capacities are invariant under biholomorphisms between bounded domains.
\end{proposition}
\begin{proof}
Monotonicity and subadditivity follow from these properties for each positive measure, then taking the supremum. For outer capacity, choose open covers with errors $\varepsilon2^{-j}$ and take their union. The bounded CLN estimate gives compact finiteness. Inner regularity for each locally finite Monge--Amp\`ere measure, followed by interchange of the two suprema, proves the compact approximation for Borel sets. Monotone convergence and the same interchange prove continuity on increasing Borel sequences. A rescaled positive quadratic function lies between zero and one on the bounded $D$, and its determinant measure is a fixed positive multiple of Lebesgue measure, giving the volume bound.

Restriction of a competitor on $D$ proves the first domain inequality. For the reverse, write $B=B(a,r)$ and choose $A>0$ so that $\rho=A(|z-a|^2-r^2)<-1$ near $K$. For a competitor $v$ on $B$, paste $\max\{v-1,\rho\}$ inside $B$ to $\rho$ outside. The boundary gluing theorem \cref{prop:psh-max-gluing} applies. The pasted function $w$ equals $v-1$ near $K$, is at least $-1$ throughout $D$, and is at most $\Lambda=\max\{0,\sup_D\rho\}$, independently of $v$. The competitor $(w+1)/(\Lambda+1)$ on $D$ and Euclidean locality give the second inequality with $C=(\Lambda+1)^n$. For outer capacities first enlarge $K$ to a compact neighborhood still inside $B$, restrict open covers to its interior, and apply these Borel inequalities there before taking their infima.

A biholomorphism bijects the capacity competitors by composition, by \cref{prop:psh-holomorphic-pullback}. Coordinate invariance of the recursive product identifies their top-degree measures by pullback. Changing variables in their integrals gives equality of the Borel capacities; the homeomorphism also bijects open covers, giving equality of outer capacities.
\end{proof}

\begin{proposition}\label{prop:bt-poles-zero-capacity}
For $K\Subset L^\circ\subset L\Subset D$ and $V\in\PSH(D)$,
\[
 \operatorname{Cap}_D^*(K\cap\{V<-s\})
                 \le\frac{C_{K,L}}{s}\|V\|_{L^1(L)},\qquad s>0.
\]
Consequently a locally contained plurisubharmonic pole set has zero outer capacity on compact subsets of any containing bounded domain. Each bounded-potential Monge--Amp\`ere measure gives that set zero mass.
\end{proposition}
\begin{proof}
Choose an open neighborhood $U$ of $K$ with compact closure in $L^\circ$. For every capacity competitor $v$, \cref{thm:cln-weighted} gives
$\int_{U\cap\{V<-s\}}(\ddc v)^n\le s^{-1}\int_{\overline U}|V|(\ddc v)^n
\le C s^{-1}\|V\|_{L^1(L)}$.
The sublevel set intersected with $U$ is open, so this bounds its outer capacity as well and proves the displayed estimate. Let $s$ increase to infinity. For local pole sets, a countable cover by compact neighborhoods of their local defining functions, domain restriction, and outer subadditivity give the result. The weighted CLN theorem also shows directly that a bounded-product measure cannot charge a pole where its defining $V$ would fail to be integrable.
\end{proof}

\begin{proposition}[B\l{}ocki]\label{prop:blocki-difference-estimate}
Let $u,h$ be bounded plurisubharmonic near the closure of a ball $B$, with $u\le h$ and $u=h$ on a collar of $\partial B$. For a bounded plurisubharmonic $v$ near $\overline B$, put $M=\sup_Bv-\inf_Bv$. Then
\[
 \int_B(h-u)^{n+1}(\ddc v)^n
       \le(n+1)!M^n\int_B(h-u)(\ddc u)^n.
\]
\end{proposition}
\begin{proof}
First suppose the functions are smooth, subtract $\sup_Bv$ from $v$, and write $w=h-u\ge0$, which vanishes on a collar. For $p\ge2$ and every closed positive complementary product $R$, the chain rule gives
\[
 \ddc(w^p)=p w^{p-1}(\ddc h-\ddc u)
               +p(p-1)w^{p-2}\mathrm{d}w\wedge\mathrm{d}^{\mathrm c}w
 \ge-pw^{p-1}\ddc u.
\]
The gradient form is $(\mathrm{i}/2\pi)\partial w\wedge\bar\partial w$, hence positive. Since $-M\le v\le0$, compact-support integration by parts yields
\[
 \int_Bw^p\ddc v\wedge R
   =\int_Bv\ddc(w^p)\wedge R
   \le pM\int_Bw^{p-1}\ddc u\wedge R.
\]
Start with $p=n+1$ and $R=(\ddc v)^{n-1}$. Repeat exactly $n$ times, ending with $p=2$, to obtain the displayed inequality.

In the bounded case choose a smaller concentric ball $B'$ whose boundary lies inside the common collar and which contains the support of $h-u$. Convolve $u,h,v$ with the same sufficiently small radial kernel. Their order is preserved, and the smoothed $h-u$ still vanishes on a fixed collar of $\partial B'$. The smoothing of $v$ on $B'$ uses only values in $B$, so its oscillation there is at most $M$. Apply the smooth inequality on $B'$ and pass to the limit by \cref{cor:bt-polynomial-weights}; all weighted differences have compact support in the fixed interior. The limiting integrands vanish on $B\setminus B'$, giving the stated inequality on $B$ with the power $M^n$.
\end{proof}
\index{B\l{}ocki's estimate}

\section{Quasi-continuity and convergence in capacity}
\label{sec:bt-quasicontinuity}

Let $D\subset\mathbb{C}^n$ remain bounded and open. A real Borel function is quasi-continuous for $\operatorname{Cap}_D$ if, for every compact $K\Subset D$ and $\varepsilon>0$, an open $G\subset D$ with $\operatorname{Cap}_D(G)<\varepsilon$ makes its restriction to $K\setminus G$ continuous. We use the same definition for extended-real functions, with the interval topology on $[-\infty,\infty]$.
\index{quasi-continuity!Monge--Amp\`ere capacity}

\begin{theorem}\label{thm:psh-capacity-quasicontinuity}
Every plurisubharmonic function on $D$ is quasi-continuous. The exceptional open set can be chosen to contain its pole set, so that the remaining function is finite. Given $\varepsilon>0$, one can also choose a single open $G$ of capacity less than $\varepsilon$ on whose complement the function is continuous throughout $D$.
\end{theorem}
\begin{proof}
First suppose $u$ is locally bounded. Work near a compact subset of a ball with closure in $D$, normalize $u$ on a larger ball, and use the common-collar modification from \cref{thm:bt-decreasing-continuity}. It gives a bounded $g$ equal to $u$ near the compact set, and continuous plurisubharmonic $g_j\downarrow g$ with one common quadratic collar. Namely take the maximum of a smooth radial approximation of $u$ and that quadratic function. For a capacity competitor $v$ on $D$, \cref{prop:blocki-difference-estimate} and Chebyshev's inequality give
\[
 \operatorname{Cap}_D(\{g_j-g>\delta\})
 \le\frac{(n+1)!}{\delta^{n+1}}
       \int_B(g_j-g)(\ddc g)^n\longrightarrow0.
\]
Here the set is understood inside $B$ and extended by the empty set outside. It is open, since $g_j$ is continuous and $g$ upper semicontinuous, and it stays away from the boundary by the common collar. The fixed measure $(\ddc g)^n$ has finite mass on $\overline B$. Bounded convergence proves the limit, and taking the supremum over $v$ preserves the displayed bound.

Choose increasing indices $j_l$ so that the open sets $G_l=\{g_{j_l}-g>2^{-l}\}$ have capacities below $\varepsilon2^{-l-1}$. Outside their union the continuous $g_{j_l}$ converge uniformly to $g$. Its restriction there is continuous. A finite ball cover and a division of the capacity budget prove the bounded case on any compact set.

For an arbitrary $u$, \cref{prop:bt-poles-zero-capacity} gives an open sublevel neighborhood $G_1$ of capacity below $\varepsilon/2$ containing all poles on the chosen compact neighborhood. Choose its level $-s$ and apply the bounded case to $\max\{u,-s\}$, with budget $\varepsilon/2$. Outside the union of the two exceptional sets this truncation equals $u$ and is finite and continuous. Finally exhaust $D$ by compact neighborhoods $K_l$ with $K_l\subset K_{l+1}^\circ$ and distribute a total budget below $\varepsilon$ over them. The union of the resulting open sets has the stated capacity. Each point of its complement lies in the interior of some $K_l$, where the relative continuity has been proved. The union also contains every pole.
\end{proof}

We need a continuous extension only from a compact set. If $f$ is continuous on a compact $F\subset\mathbb{R}^d$, it has a bounded continuous extension to $\mathbb{R}^d$ with the same bound. Here is a direct construction. If $|f|\le a$, the clamped functions
$\max\{-a,\min\{a,\inf_{y\in F}(f(y)+L|x-y|)\}\}$
are Lipschitz and converge uniformly to $f$ on $F$ as $L\to\infty$. Uniform continuity proves this: near $x\in F$ the values of $f$ differ little, while points at a fixed positive distance are excluded from the infimum when $L$ is large. Choose such an approximation with error at most $a/2$, repeat for its residual on $F$, and sum the bounded extensions of these residual approximations. The bounds form a geometric series, so the sum converges uniformly and equals $f$ on $F$. Clamping that sum to $[-a,a]$ preserves its values on $F$ and gives the desired bound. The empty $F$ and $a=0$ use the zero extension.

\begin{lemma}\label{lem:quasi-continuous-level-capacity}
If $L\Subset D$ is compact, $f$ is bounded and quasi-continuous on a neighborhood of $L$, and $K\Subset L^\circ$, then, for $\delta>0$,
\[
 \operatorname{Cap}_D^*(K\cap\{|f|>\delta\})
       \le\operatorname{Cap}_D(L\cap\{|f|>\delta/2\}).
\]
\end{lemma}
\begin{proof}
Choose open $G$ of capacity below $\eta$ making $f$ continuous on $L\setminus G$, and extend that restriction continuously with the same bounded norm by the preceding construction. Write the extension as $g$. If $K\subset U\Subset L^\circ$, the open set
$G\cup(U\cap\{|g|>\delta/2\})$ covers the set on the left. Its portion outside $G$ lies in $L\cap\{|f|>\delta/2\}$. Subadditivity bounds its capacity by the right side plus $2\eta$. Let $\eta\downarrow0$.
\end{proof}

\begin{definition}\label{def:local-capacity-convergence}
A net of Borel functions $f_a:D\to[-\infty,+\infty]$ converges locally in capacity to a Borel function $f$ if, for every compact $K\Subset D$ and $\delta>0$,
\[
 \operatorname{Cap}_D^*(K\cap\{f_a>f+\delta\})\longrightarrow0,
 \qquad
 \operatorname{Cap}_D^*(K\cap\{f_a<f-\delta\})\longrightarrow0.
\]
Addition of a finite constant preserves either infinite value in these ordered comparisons. For finite-valued functions the two conditions are equivalent to $\operatorname{Cap}_D^*(K\cap\{|f_a-f|>\delta\})\to0$, by subadditivity. Local comparisons in \cref{prop:bt-capacity-basic} make the definition independent of the containing coordinate ball on compact subsets.
\end{definition}
\index{convergence!Monge--Amp\`ere capacity}

\begin{proposition}\label{prop:psh-decreasing-capacity}
If $u_j\downarrow u$ are locally bounded plurisubharmonic functions with locally bounded limit, then $u_j\to u$ locally in capacity. Locally uniformly bounded functions converging in capacity also converge in $L^1_{\loc}$.
\end{proposition}
\begin{proof}
Modify the decreasing functions on a ball to have the common collar and agree with their original values on a compact neighborhood of the test set. Put $w_j=u_j-u\ge0$ there. Each $w_j$ is bounded and quasi-continuous, since it is a difference of two such functions. The B\l{}ocki estimate bounds the inner capacity of its positive superlevel sets by
$C\delta^{-n-1}\int_B w_j(\ddc u)^n$, which tends to zero by bounded convergence. \Cref{lem:quasi-continuous-level-capacity} turns this into the required outer-capacity bound on the smaller compact set, with threshold $\delta/2$. A finite cover proves the local assertion.

For the $L^1$ implication, if $|f_j|,|f|\le M$ near a compact $K$, split the integral into $\{|f_j-f|>\delta\}$ and its complement. The volume bound from \cref{prop:bt-capacity-basic} gives
\[
 \int_K|f_j-f|\,\mathrm{d}V
 \le 2MC_K\operatorname{Cap}_D^*(K\cap\{|f_j-f|>\delta\})
          +\delta|K|.
\]
Let $j$ increase and then $\delta$ decrease.
\end{proof}

\begin{lemma}\label{lem:bounded-product-capacity-domination}
On a ball $B$, the trace measures of products of locally uniformly bounded plurisubharmonic potentials are uniformly dominated on smaller compact sets by a constant times $\operatorname{Cap}_B$. Their coefficient variations have the same bound. If such measures $\mu_j\to\mu$ weakly, then $f_j\mu_j\to f\mu$ weakly whenever the real Borel $f_j$ are locally uniformly bounded, converge in capacity to a bounded quasi-continuous $f$, and all products have the indicated common bounds. This includes a fixed quasi-continuous weight $f_j=f$.
\end{lemma}
\begin{proof}
Restrict to a slightly smaller ball on which every potential has the same finite bound and normalize them to $0\le v_l\le1$. Also normalize the quadratic function $\phi=|z-a|^2/R^2$ on that ball. The function
$F=(v_1+\cdots+v_q+(n-q)\phi)/n$ is a capacity competitor. Positivity and the symmetric expansion of $(\ddc F)^n$ show that it dominates a fixed positive multiple of
$\ddc v_1\wedge\cdots\wedge\ddc v_q\wedge\beta^{n-q}$: the expansion contains all permutations of those $n$ nonnegative factors. Undoing the normalizations gives the trace bound. The current coefficient bound follows from \cref{thm:positive-current-order-zero}. Domain comparison transfers the estimate back to the chosen $\operatorname{Cap}_B$ on smaller compact sets.

For a fixed bounded quasi-continuous $f$, choose an open set $G$ of arbitrarily small capacity making it continuous on a compact neighborhood of the test support. Extend that restriction continuously with the same bound, as above. The extended weight converges against $\mu_j$ by weak convergence; continuous compact tests are allowed by uniform mass bounds and smooth uniform approximation. The errors are confined to $G$ and bounded by its capacity times a fixed constant. The limit measure has the same outer-capacity domination on the interior: for open subsets this follows from the vague set inequality in \cref{prop:vague-set-inequalities}, and arbitrary subsets are bounded by open covers. Let the capacity of $G$ decrease to zero.

For varying weights, split $|f_j-f|$ at level $\delta$ on a fixed compact neighborhood. The integral of the difference against $\mu_j$ is at most $\delta\mu_j(K)$ plus a common bound times the outer capacity of its superlevel set. Capacity convergence makes the second term tend to zero. Then let $\delta\downarrow0$ and apply the fixed-weight result. The same argument for coefficient variations proves the corresponding current statement.
\end{proof}

\begin{theorem}\label{thm:bt-capacity-continuity}
If $u_{l,j}\to u_l$ locally in capacity for locally uniformly bounded plurisubharmonic functions with locally bounded plurisubharmonic limits, then
\[
 \ddc u_{1,j}\wedge\cdots\wedge\ddc u_{q,j}
             \longrightarrow\ddc u_1\wedge\cdots\wedge\ddc u_q
\]
weakly, for $q\le n$. One may also insert locally uniformly bounded Borel weights $f_j\to f$ in capacity when $f$ is quasi-continuous. The base current in this assertion is $1$.
\end{theorem}
\begin{proof}
Capacity convergence gives $L^1$ convergence by \cref{prop:psh-decreasing-capacity}, so the one-factor Hessians converge distributionally. For the induction step, the $(q-1)$-factor products converge weakly and their coefficient variations have uniform capacity bounds by \cref{lem:bounded-product-capacity-domination}. That lemma passes multiplication by $u_{1,j}$ to multiplication by $u_1$. Apply $\ddc$ to obtain the $q$-factor assertion. The same lemma then proves the additional-weight assertion. All estimates are local on a bounded ball; the domain comparison makes the capacity hypotheses valid there.
\end{proof}

\section{Increasing limits and exceptional envelopes}
\label{sec:bt-increasing}

Increasing potentials can have a pointwise supremum below its canonical regularization on an exceptional set. Quasi-continuity controls fixed weights against varying product measures, allowing the increasing-limit argument to identify that canonical value before the exceptional set is analyzed.

\begin{theorem}\label{thm:bt-increasing-continuity}
Let $u_{l,j}$ be locally uniformly bounded plurisubharmonic functions increasing almost everywhere to locally bounded plurisubharmonic $u_l$, for $0\le l\le q\le n$. Then
\[
 \begin{aligned}
 \ddc u_{1,j}\wedge\cdots\wedge\ddc u_{q,j}
 &\longrightarrow\ddc u_1\wedge\cdots\wedge\ddc u_q,\\
 u_{0,j}\ddc u_{1,j}\wedge\cdots\wedge\ddc u_{q,j}
 &\longrightarrow u_0\ddc u_1\wedge\cdots\wedge\ddc u_q
 \end{aligned}
\]
weakly. Moreover $u_{l,j}\to u_l$ locally in capacity.
\end{theorem}
\begin{proof}
Induct on the weighted assertion with $q$ Hessian factors. For $q=0$, bounded convergence with respect to Lebesgue measure applies. The weighted assertion for $q-1$, followed by $\ddc$, proves the unweighted assertion for $q$. Denote its products by $S_j\to S$. As in the decreasing proof, work on a ball with a common quadratic collar, leaving the functions unchanged near the smaller test set. The maximum modification preserves the increasing sequences and their almost-everywhere limits. Local bounds give weakly convergent subsequences of $u_{0,j}S_j$; call any limit $\Theta$.

Every $u_{0,j}\le u_0$ pointwise. Indeed the inequality holds almost everywhere, and uniqueness of the canonical representative of their maximum makes it pointwise. Continuous upper majorants of $u_0$, followed by their decreasing limit as in \cref{lem:usc-weight-weak-limit}, give $\Theta\le u_0S$. The difference is positive and supported away from the common collar. Write $\gamma=\beta^{n-q}$.

For a fixed $k$ and all $j\ge k$, positivity gives
$\int_Bu_{0,j}S_j\wedge\gamma\ge\int_Bu_{0,k}S_j\wedge\gamma$.
The fixed $u_{0,k}$ is quasi-continuous; \cref{lem:bounded-product-capacity-domination} passes the right side to $\int_Bu_{0,k}S\wedge\gamma$. The bounded common-collar identity and symmetry rewrite this as
\[
 \int_Bu_1\ddc u_{0,k}\wedge\ddc u_2\wedge\cdots\wedge\ddc u_q\wedge\gamma.
\]
The unweighted $q$-factor increasing assertion, already obtained from the induction hypothesis, applies to the product in this display as $k\to\infty$. The fixed quasi-continuous $u_1$ can again be passed to its limit. A final common-collar interchange identifies that limit with $\int_Bu_0S\wedge\gamma$. Collar integrals converge separately as in \cref{cor:bt-polynomial-weights}; thus these statements about integrals use valid compact tests. We have proved that the trace mass of $\Theta$ is at least that of $u_0S$. Positivity of their difference gives equality and then $\Theta=u_0S$. This proves weighted continuity and completes the induction.

Having proved both continuity assertions for every $q\le n$, use the same collar modification for capacity convergence, with $w_j=u_l-u_{l,j}\ge0$. The B\l{}ocki estimate now bounds its superlevel capacity by a constant times
\[
 \int_B(u_l-u_{l,j})(\ddc u_{l,j})^n\longrightarrow0.
\]
The term with fixed weight $u_l$ converges by quasi-continuity and the unweighted assertion just proved; the term with weight $u_{l,j}$ converges by its weighted assertion. Their limits cancel. \Cref{lem:quasi-continuous-level-capacity} gives outer-capacity convergence on smaller compact sets. Undo the collar modification and use a finite cover.
\end{proof}
\index{Bedford--Taylor product!increasing continuity}

\begin{corollary}\label{cor:psh-envelope-zero-capacity}
For a nonempty locally uniformly bounded-above family of plurisubharmonic functions on $D$, the set where its pointwise supremum differs from its upper semicontinuous regularization has zero outer Monge--Amp\`ere capacity on compact subsets of $D$.
\end{corollary}
\begin{proof}
Choose a countable family with the same regularized envelope by \cref{thm:psh-upper-envelopes}, and take its increasing finite maxima $v_j$. They converge almost everywhere to the plurisubharmonic regularization $u$. On a fixed relatively compact coordinate neighborhood, the functions $\max\{v_j,-m\}$ are locally uniformly bounded and increase almost everywhere to $\max\{u,-m\}$. By \cref{thm:bt-increasing-continuity}, their convergence is in capacity. If their pointwise supremum is below the latter limit, it is below by some positive rational amount. Every such discrepancy set is contained in every sufficiently late superlevel difference set, whose outer capacities tend to zero. Countable subadditivity therefore gives zero outer capacity for that bounded discrepancy.

Any point where $\sup_jv_j<u$ has $u>-m$ for some integer $m$, and hence lies in one of these truncated discrepancies. Another countable union proves the assertion for the sequence. The full pointwise family supremum dominates $\sup_jv_j$ and is at most $u$, so its discrepancy is a subset of the same zero-capacity set. Outer capacity is defined for arbitrary subsets, avoiding a Borel assumption on an unrestricted supremum.
\end{proof}

\index{plurisubharmonic function!capacity-zero envelope discrepancy}

\section{The homogeneous Dirichlet problem and balayage}
\label{sec:bt-homogeneous-dirichlet}

An extremal function is maximal where its obstacle is absent. To translate that comparison property into vanishing Monge--Amp\`ere measure, we need a local replacement solving the homogeneous equation. Ball automorphisms give the second-difference bound which makes the smooth-boundary envelope twice differentiable almost everywhere.

\begin{lemma}\label{lem:ball-centered-automorphisms}
For $a$ in the unit ball $B\subset\mathbb C^n$ there are automorphisms $S_{a,h}$ of $B$, defined for small $h$, with $S_{a,0}$ the identity and $S_{a,h}(a)=a+h$. They extend smoothly to $\overline B$. On each compact set of centers $a$, their first two derivatives in the real parameter $h$ are uniformly bounded on $\overline B$.
\end{lemma}
\begin{proof}
Use the Hermitian product $\langle z,a\rangle=\sum z_j\bar a_j$, put $s_a=\sqrt{1-|a|^2}$, and define
\[
 T_a(z)=\frac{s_a z+\langle z,a\rangle a/(1+s_a)-a}
                       {1-\langle z,a\rangle}.
\]
This formula is smooth at $a=0$, where $T_0$ is the identity. In unitary coordinates with $a=(t,0,\ldots,0)$, $t\ge0$, it reads
\[
 T_a(z)=\left(\frac{z_1-t}{1-tz_1},
           \frac{\sqrt{1-t^2}\,z_2}{1-tz_1},\ldots,
           \frac{\sqrt{1-t^2}\,z_n}{1-tz_1}\right).
\]
Direct substitution gives $T_a(a)=0$, $T_{-a}\circ T_a=\mathrm{id}$, and
\[
 1-|T_a(z)|^2=\frac{(1-|a|^2)(1-|z|^2)}{|1-\langle z,a\rangle|^2}.
\]
Thus $T_a$ bijects $B$ and preserves its sphere. Set $S_{a,h}=T_{-(a+h)}\circ T_a$. It sends $a$ to $a+h$ and equals the identity for $h=0$. If $a$ ranges in a compact subset and $h$ is sufficiently small, the denominators in both maps stay uniformly away from zero on $\overline B$. The smooth formula gives the stated derivative bounds.
\end{proof}

\begin{theorem}\label{thm:homogeneous-ball-dirichlet}
For a ball $B$ and a real continuous $\varphi$ on $\partial B$, there is a unique $H\in\PSH(B)\cap C(\overline B)$ with
\[
 (\ddc H)^n=0\quad\hbox{on }B,\qquad H|_{\partial B}=\varphi.
\]
It is the envelope of all plurisubharmonic $v$ with upper boundary limits at most $\varphi$. If $\varphi$ is the restriction of a $C^2$ function on a neighborhood of the sphere, then $H\in C^{1,1}_{\loc}(B)$.
\end{theorem}
\begin{proof}
Scale to the unit ball. First suppose $\varphi$ is smooth. Extend it smoothly to $\overline B$, for instance radially on an annular collar and then with a smooth cutoff. For sufficiently large $A$, the extension plus $A(|z|^2-1)$ is plurisubharmonic; call it $s$. It has boundary value $\varphi$. The ordinary harmonic solution $P\varphi$ from \cref{thm:ball-dirichlet} bounds every competitor above, by the subharmonic comparison principle. Thus the envelope $H$ satisfies $s\le H\le P\varphi$. Its upper regularization is plurisubharmonic by \cref{thm:psh-upper-envelopes}. The same barriers show that this regularization has boundary value $\varphi$, so it is a competitor itself. It therefore equals $H$ pointwise. In particular $H$ approaches $\varphi$ continuously at the sphere and is bounded.

For $a$ in a fixed compact subset, apply \cref{lem:ball-centered-automorphisms}. On the sphere, Taylor's formula in $h$ gives
\[
 \tfrac12\bigl(\varphi\circ S_{a,h}+\varphi\circ S_{a,-h}\bigr)
                 \le\varphi+C_K|h|^2.
\]
The constant is uniform in $a$ and in the boundary point. The average of $H\circ S_{a,h}$ and $H\circ S_{a,-h}$, minus $C_K|h|^2$, belongs to the envelope class. Evaluating its inequality with $H$ at $a$ yields
\[
 H(a+h)+H(a-h)-2H(a)\le2C_K|h|^2.
\]
Convolve on a slightly larger compact neighborhood. The same inequality holds for $H_\varepsilon$, and division by the squared step followed by its limit gives $D^2H_\varepsilon[h,h]\le2C_K|h|^2$. Plurisubharmonicity gives
\[
 D^2H_\varepsilon[h,h]+D^2H_\varepsilon[\mathrm ih,\mathrm ih]\ge0,
\]
so the real Hessians are also bounded below by $-2C_K$. The proof of \cref{prop:bounded-hessian-taylor} applies even before continuity of $H$ has been shown: the radial smoothings converge in $L^1$ and decrease pointwise to $H$, their bounded gradients give a continuous subsequential limit, and uniqueness of the canonical representative identifies it with $H$. Consequently $H$ is locally $C^{1,1}$ and has the second-order expansion of that proposition almost everywhere.

Its complex Hessian has bounded measurable coefficients. At their common Lebesgue points the smoothed Hessians tend to those coefficients, by \cref{thm:lebesgue-differentiation}; they remain bounded on compact sets. Determinant expansion, bounded convergence and \cref{prop:bt-uniform-continuity} therefore give
\[
 (\ddc H)^n=\frac{n!}{\pi^n}\det(H_{j\bar k})\,\mathrm dV_{2n}.
\]
If that determinant were positive at a point $a$ having the second-order expansion, the real quadratic Taylor polynomial $Q$ would have positive definite complex Hessian. Choose $\eta>0$ so small that $Q-\eta|z-a|^2$ is still plurisubharmonic. For a sufficiently small radius $r$, Taylor's remainder on $|z-a|=r$ has absolute value below $\eta r^2/2$. The polynomial
$Q-\eta|z-a|^2+\eta r^2/4$ is then strictly below $H$ on that sphere and strictly above it at $a$. Paste its maximum with $H$ in the smaller ball. \Cref{prop:psh-max-gluing} produces a competitor with the same boundary values which exceeds the envelope at $a$, a contradiction. Hence the determinant vanishes almost everywhere, proving the equation.

For continuous $\varphi$, extend it continuously to an annular neighborhood by radial projection and smooth there. This gives smooth boundary functions $\varphi_j\to\varphi$ uniformly. Adding a constant to the boundary data adds that constant to its envelope. Order of the classes consequently gives
$\|H_j-H_k\|_{L^\infty(B)}\le\|\varphi_j-\varphi_k\|_{L^\infty(\partial B)}$.
The envelopes converge uniformly on $\overline B$ to a continuous plurisubharmonic $H$ with boundary value $\varphi$. Uniform continuity of the bounded products gives $(\ddc H)^n=0$. The same constant-shift bounds identify it with the envelope for $\varphi$. The argument for smooth data used only two derivatives, so it also gives the stated $C^2$-boundary regularity. Finally \cref{thm:bt-comparison}, applied in both directions to any two solutions, proves uniqueness.
\end{proof}
\index{Dirichlet problem!homogeneous Monge--Amp\`ere}

\begin{corollary}\label{cor:homogeneous-psh-balayage}
Let $u$ be locally bounded and plurisubharmonic on $\Omega$, and let $B\Subset\Omega$ be a ball. There is a locally bounded plurisubharmonic $\widehat u$ on $\Omega$ which equals $u$ outside $B$, is at least $u$ inside $B$, and satisfies $(\ddc\widehat u)^n=0$ in $B$. The construction is increasing in $u$, and its restriction to $B$ is the largest plurisubharmonic function with upper boundary limits at most $u|_{\partial B}$.
\end{corollary}
\begin{proof}
Take smooth radial regularizations $u_j\downarrow u$ on a neighborhood of $\overline B$ and solve the homogeneous problem with their continuous boundary values. Write its solutions as $H_j$. Comparison gives $u\le u_j\le H_j$ in $B$, and order of the boundary envelopes gives $H_j\downarrow H$. A common bound for $u$ and $u_j$ near $\overline B$ bounds this sequence above and below, so $H$ is bounded and plurisubharmonic; decreasing continuity gives $(\ddc H)^n=0$. For $x\in\partial B$,
\[
 \limsup_{B\ni z\to x}H(z)\le u_j(x)\quad\hbox{for every }j,
 \qquad\limsup_{B\ni z\to x}H(z)\le u(x).
\]
Paste $H$ inside $B$ to $u$ outside using \cref{prop:psh-max-gluing}; this is $\widehat u$.

Every competitor $v$ with upper boundary limits at most $u|_{\partial B}$ is dominated by each $H_j$. To apply bounded comparison, first truncate $v$ below a constant smaller than the minimum of $H_j$ on $\overline B$; its upper bound follows from the subharmonic maximum principle. The boundary upper limits of this truncation are still at most $u_j$. Comparison with the continuous zero-measure $H_j$ gives the required domination. Thus $v\le H$. Conversely $H$ belongs to this class by the displayed upper boundary estimate, so it is its largest member. Inclusion of boundary classes gives monotonicity in $u$. This proves the asserted properties of the canonical replacement.
\end{proof}
\index{balayage!plurisubharmonic homogeneous replacement}

\begin{definition}\label{def:maximal-plurisubharmonic}
A plurisubharmonic $u$ on $\Omega$ is maximal if every plurisubharmonic $v$ with $v\le u$ outside a compact subset of $\Omega$ satisfies $v\le u$ everywhere.
\end{definition}
\index{plurisubharmonic function!maximal}

\begin{theorem}\label{thm:bounded-psh-maximality}
A locally bounded plurisubharmonic function is maximal if and only if its Monge--Amp\`ere measure vanishes. In particular maximality is local within this class.
\end{theorem}
\begin{proof}
If the measure vanishes and $v\le u$ outside a compact $K$, choose $K\Subset D\Subset\Omega$. A finite upper bound for $v$ on $\overline D$ follows from upper semicontinuity. Choose $-M$ below the lower bound of $u$ there. The bounded plurisubharmonic $\max\{v,-M\}$ is at most $u$ near $\partial D$. \Cref{thm:bt-comparison} gives the inequality on $D$, and hence $v\le u$ everywhere.

Conversely apply \cref{cor:homogeneous-psh-balayage} on any ball $B\Subset\Omega$. Its replacement dominates $u$ and equals it off the compact $\overline B$. Maximality makes the two functions equal. The measure of $u$ consequently vanishes on $B$, and hence throughout $\Omega$. Vanishing of a measure is a local property, which proves the last assertion.
\end{proof}

\section{Density lower bounds and mixed products}
\label{sec:bt-density-lower-bounds}

Let $\Omega\subset\mathbb C^n$ be open, with $n\ge1$, and put $\eta=\ddc|z|^2=\beta/\pi$. Thus $\eta^n=n!\pi^{-n}\,\mathrm dV_{2n}$ by \eqref{eq:bt-determinant-normalization}. A lower bound for a Monge--Amp\`ere density can be tested against smooth functions touching the potential from above. A matrix formula then turns these tests into linear differential inequalities, which survive convolution. This permits the smooth mixed determinant inequality to pass to locally bounded potentials.

Here a real $C^2$ function $q$ touches a real function $u$ from above at $x$ if $q(x)=u(x)$ and $q\ge u$ on a neighborhood of $x$.

\begin{lemma}\label{lem:determinant-linearization-mixed}
For a positive semidefinite Hermitian $n\times n$ matrix $A$,
\[
 (\det A)^{1/n}=\inf_{B>0,\,\det B=1}\frac{\operatorname{tr}(BA)}n.
\]
In the infimum $B>0$ means positive definite Hermitian. For positive semidefinite $A,C$, define $M_k(A,C)$ by
\[
 \det(tA+sC)=\sum_{k=0}^n\binom nk M_k(A,C)t^ks^{n-k}.
\]
For each integer $1\le k\le n-1$,
\[
 M_k(A,C)\ge(\det A)^{k/n}(\det C)^{(n-k)/n}.
\]
The endpoint coefficients are $M_n(A,C)=\det A$ and $M_0(A,C)=\det C$.
\end{lemma}
\begin{proof}
If $A>0$, apply the arithmetic--geometric mean inequality to the eigenvalues of $B^{1/2}AB^{1/2}$. Equality is attained at $B=(\det A)^{1/n}A^{-1}$. If $A$ is singular, use this choice for $A+\varepsilon I$; its trace against $A$ is at most $n\det(A+\varepsilon I)^{1/n}$ and tends to zero. All the traces are nonnegative, proving the first assertion.

Suppose first that $A,C>0$ and let $\lambda_1,\ldots,\lambda_n$ be the eigenvalues of $A^{-1/2}CA^{-1/2}$. Coefficient comparison gives
\[
 M_k(A,C)=\frac{\det A}{\binom nk}
 \sum_{|J|=n-k}\prod_{j\in J}\lambda_j.
\]
Each eigenvalue occurs in $\binom{n-1}{n-k-1}$ of these products. Their geometric mean is therefore $(\lambda_1\cdots\lambda_n)^{(n-k)/n}$. The arithmetic--geometric mean inequality gives the stated bound since $\prod_j\lambda_j=\det C/\det A$. Replace both matrices by their sums with $\varepsilon I$ and let $\varepsilon\downarrow0$ to include the semidefinite case.
\end{proof}

\begin{lemma}\label{lem:constant-elliptic-upper-tests}
Let $d\ge1$ be an integer and $L=\sum_{a,b=1}^d a_{ab}\partial_a\partial_b$, where $(a_{ab})$ is a real symmetric positive definite constant matrix. Let $w$ be real, locally bounded and upper semicontinuous on an open $U\subset\mathbb R^d$, and let $F:U\to\mathbb R$ be continuous. Suppose every $C^2$ function $q$ touching $w$ from above at $x$ satisfies $Lq(x)\ge F(x)$. Then $Lw\ge F$ distributionally.
\end{lemma}
\begin{proof}
A real linear change of variables makes $L$ the ordinary Laplacian. If $F=0$, the upper-test condition implies harmonic comparison on every relatively compact ball $B(y,r)$. Indeed, let $h$ be harmonic on that ball, continuous on its closure, and at least $w$ on the boundary. If $w-h$ is positive somewhere inside, choose $\varepsilon>0$ so small that
$w-h-\varepsilon(r^2-|x-y|^2)$ still has a positive value there. Its positive maximum is attained inside, since its boundary upper limits are nonpositive. Adding this maximum to $h+\varepsilon(r^2-|x-y|^2)$ gives an upper test with strictly negative Laplacian, a contradiction. Decreasing continuous majorants of $w|_{\partial B}$ and their Poisson extensions gives the sphere submean inequality, as in \cref{thm:subharmonic-comparison}. Consequently \cref{thm:subharmonic-distribution} gives $\Delta w\ge0$.

For general $F$, work on a smaller ball and choose smooth $F_j\le F$ converging uniformly to $F$ there. Such functions result from convolution on a larger neighborhood followed by subtraction of its uniform approximation error. Multiply $F_j$ by a smooth cutoff equal to one near the smaller ball. Convolution with the fundamental solution in \cref{prop:laplacian-fundamental-solution} gives a smooth $h_j$ with $\Delta h_j=F_j$ on that ball: derivatives may be placed on the compactly supported smooth factor. Every upper test for $w-h_j$ has nonnegative Laplacian, so the first case gives $\Delta w\ge F_j$ distributionally. Uniform convergence gives $\Delta w\ge F$. Returning to the original coordinates proves the lemma.
\end{proof}

For a real $C^2$ function $q$ write $H(q)=(q_{j\bar k})$. For a positive definite Hermitian $B$ with $\det B=1$, put $L_Bq=\operatorname{tr}(BH(q))$. This is a constant real elliptic operator: diagonalizing $B$ by a unitary coordinate change expresses it as a positive weighted sum of the planar Laplacians, each with factor $1/4$.

\begin{proposition}\label{prop:bounded-monge-ampere-linear-density}
Let $u$ be locally bounded and plurisubharmonic on $\Omega$, and let $f\ge0$ be continuous. The following conditions are equivalent:
\begin{enumerate}
\item $(\ddc u)^n\ge f\eta^n$ as measures;
\item every real $C^2$ function $q$ touching $u$ from above at $x$ satisfies $\det H(q)(x)\ge f(x)$;
\item $L_Bu\ge n f^{1/n}$ distributionally for every positive definite Hermitian $B$ with $\det B=1$.
\end{enumerate}
\end{proposition}
\begin{proof}
The complex Hessian of any upper test $q$ is positive semidefinite at the contact point. Restrict to a complex line through that point and use the submean inequality for $u$; the circle mean of $q$ is at least its center value. The second-order expansion gives nonnegativity of $H(q)$ on every complex direction.

Assume (1), and normalize a test by $q(x)=u(x)$ and $q\ge u$ near $x$. If $\det H(q)(x)<f(x)$, choose $\delta>0$ so small that $q_\delta=q+\delta|z-x|^2$ has positive definite complex Hessian at $x$ and still has determinant below $f(x)$. On a sufficiently small closed ball $\overline{B(x,r)}$, it is strictly plurisubharmonic and $\det H(q_\delta)<f$. Set $b=q_\delta-\delta r^2/2$. On the sphere $b>u$, whereas $b(x)<u(x)$. The Borel set $S=\{b<u\}\cap B(x,r)$ has positive Lebesgue measure. Otherwise $u\le b$ almost everywhere on the ball, and uniqueness of the plurisubharmonic representative in \cref{thm:psh-distribution-representative}, applied to $\max\{u,b\}$, would imply $u\le b$ everywhere. Bounded comparison in \cref{thm:bt-comparison} now gives
\[
 \int_S f\eta^n\le\int_S(\ddc u)^n
 \le\int_S(\ddc b)^n
 =\int_S\det H(q_\delta)\eta^n,
\]
contradicting the strict density inequality. The measures here are finite since the ball is relatively compact in the neighborhood where the potentials are bounded. Thus (2) holds.

Assume (2). The first formula of \cref{lem:determinant-linearization-mixed} gives $L_Bq(x)\ge n f(x)^{1/n}$ for every upper test and each fixed admissible $B$. Apply \cref{lem:constant-elliptic-upper-tests} with the continuous function $n f^{1/n}$ to obtain (3).

Finally assume (3). Choose the nonnegative radial regularizing kernels used in \cref{thm:psh-distribution-representative}, so that $u_\varepsilon=u*\rho_\varepsilon\downarrow u$ on smaller open sets. Convolving each distributional inequality gives
\[
 \operatorname{tr}(BH(u_\varepsilon))\ge n(f^{1/n}*\rho_\varepsilon)
 \qquad(B>0,\ \det B=1).
\]
These are pointwise inequalities between smooth functions. The matrix infimum formula therefore gives $\det H(u_\varepsilon)\ge(f^{1/n}*\rho_\varepsilon)^n$. The right side tends locally uniformly to $f$. The determinant normalization and the weak decreasing continuity of \cref{thm:bt-decreasing-continuity} give (1). This argument includes vanishing densities and singular Hessians.
\end{proof}
\index{Monge--Amp\`ere measure!density lower bound}

\begin{corollary}\label{cor:bounded-mixed-continuous-density}
Let $u,v$ be locally bounded plurisubharmonic functions on $\Omega$, and let $f,g\ge0$ be continuous with $(\ddc u)^n\ge f\eta^n$ and $(\ddc v)^n\ge g\eta^n$. For each integer $1\le k\le n-1$,
\[
 (\ddc u)^k\wedge(\ddc v)^{n-k}
 \ge f^{k/n}g^{(n-k)/n}\eta^n.
\]
For $k=0,n$ the corresponding assertions are the given pure-measure bounds.
\end{corollary}
\begin{proof}
Use the same radial kernels for both potentials. The proof of \cref{prop:bounded-monge-ampere-linear-density} gives
\[
 \det H(u_\varepsilon)\ge(f^{1/n}*\rho_\varepsilon)^n,
 \qquad
 \det H(v_\varepsilon)\ge(g^{1/n}*\rho_\varepsilon)^n.
\]
Expand $(t\ddc u_\varepsilon+s\ddc v_\varepsilon)^n$ and its smooth determinant density. The coefficient of $\binom nk t^ks^{n-k}$ is $M_k(H(u_\varepsilon),H(v_\varepsilon))\eta^n$. The second assertion of \cref{lem:determinant-linearization-mixed} bounds this below by
$(f^{1/n}*\rho_\varepsilon)^k(g^{1/n}*\rho_\varepsilon)^{n-k}\eta^n$.
The continuous density converges locally uniformly to $f^{k/n}g^{(n-k)/n}$. Mixed decreasing continuity in \cref{thm:bt-decreasing-continuity} passes the measure inequality to the limit.
\end{proof}
\index{Bedford--Taylor product!mixed density inequality}


\subsection{Several independent factors}
\label{sec:bt-multifactor-density}
\Cref{cor:bounded-mixed-continuous-density} concerns products in which only two Hessians occur, each repeated several times. For a product $\ddc u_1\wedge\cdots\wedge\ddc u_n$ of $n$ arbitrary potentials, the coefficients $M_k$ of \cref{lem:determinant-linearization-mixed} are replaced by the mixed discriminant $D(A_1,\ldots,A_n)$ of Hermitian matrices. Recall from the paragraph before \cref{lem:mixed-discriminant-basic} that $D(A_1,\ldots,A_n)$ is $1/n!$ times the coefficient of $t_1\cdots t_n$ in $\det(\sum_jt_jA_j)$; it is symmetric and real multilinear, and $D(A,\ldots,A)=\det A$ (\cref{lem:mixed-discriminant-basic}(1)). We prove the geometric-mean inequality for $D$ by means of a reverse Cauchy--Schwarz inequality, which follows from the signature of $D$ computed in \cref{lem:mixed-discriminant-signature}. Convolution then extends it to locally bounded potentials as before.

\begin{lemma}\label{lem:multifactor-mixed-discriminant}
Let $n\geq1$. For positive semidefinite Hermitian $n\times n$ matrices $A_1,\ldots,A_n$,
\[
 D(A_1,\ldots,A_n)\geq\prod_{i=1}^n(\det A_i)^{1/n}.
\]
\end{lemma}
\begin{proof}
For $n=1$ both sides equal $A_1$. Let $n\geq2$, and assume first that all $A_i$ are positive definite. By \cref{lem:mixed-discriminant-basic}(3), $D$ is positive on positive definite matrices.

We first derive the reverse Cauchy--Schwarz inequality. Fix positive definite matrices $A_3,\ldots,A_n$, and consider the symmetric bilinear form
\[
 q(A,B)=D(A,B,A_3,\ldots,A_n)
\]
on Hermitian matrices. Let $A,B$ be positive definite, put $c=q(A,B)/q(B,B)$ and $M=A-cB$. Then $q(M,B)=0$, so $q(M,M)\leq0$ by \cref{lem:mixed-discriminant-signature}. Since $q(M,M)=q(A,A)-q(A,B)^2/q(B,B)$, we obtain
\begin{equation}\label{eq:mixed-discriminant-reverse-cauchy}
 q(A,B)^2\geq q(A,A)q(B,B).
\end{equation}

Put $\widetilde A_i=(\det A_i)^{-1/n}A_i$, so that $\det\widetilde A_i=1$. By multilinearity it suffices to prove that $D(\widetilde A_1,\ldots,\widetilde A_n)\geq1$. Consider the finitely many ordered $n$-tuples whose entries belong to the set $\{\widetilde A_1,\ldots,\widetilde A_n\}$, and let $F>0$ be the smallest mixed discriminant of such a tuple. A tuple whose entries all equal $\widetilde A_i$ has discriminant $\det\widetilde A_i=1$, so $F\leq1$. The multiplicity of a matrix in a tuple is the number of entries equal to it. Among the tuples with discriminant $F$, choose one for which the sum of the squares of the multiplicities is largest.

Suppose that this tuple contains two different matrices $A$ and $B$. By symmetry of $D$ we may place them in the first two slots, so that $F=q(A,B)$, where $q$ is formed with the remaining entries of the tuple. Replacing the pair $(A,B)$ by $(A,A)$ or by $(B,B)$ gives tuples with discriminants $q(A,A)\geq F$ and $q(B,B)\geq F$, whereas \eqref{eq:mixed-discriminant-reverse-cauchy} gives $q(A,A)q(B,B)\leq F^2$. Hence $q(A,A)=q(B,B)=F$, and both new tuples are minimizing. If $A$ and $B$ have multiplicities $r$ and $s$, the two replacements change the sum of the squared multiplicities by $2(r-s)+2$ and $2(s-r)+2$, respectively. These two numbers add up to $4$, so one of them is positive, which contradicts the choice of the tuple. Therefore all entries of the chosen tuple are equal, and $F=1$. In particular $D(\widetilde A_1,\ldots,\widetilde A_n)\geq1$.

Finally, for positive semidefinite $A_i$, apply the inequality to $A_i+\varepsilon I$ and let $\varepsilon\downarrow0$. Both sides are continuous: $D$ is a polynomial in the matrix entries, and $(\det A)^{1/n}$ is continuous on positive semidefinite matrices.
\end{proof}
\index{mixed discriminant!geometric-mean inequality}

\begin{corollary}\label{cor:bounded-multifactor-continuous-density}
Let $u_1,\ldots,u_n$ be locally bounded plurisubharmonic functions on an open set $\Omega\subset\mathbb C^n$, $n\geq1$, and let $f_1,\ldots,f_n\geq0$ be continuous functions on $\Omega$ such that
\[
 (\ddc u_i)^n\geq f_i\eta^n\quad(1\leq i\leq n),
 \qquad\eta=\ddc|z|^2.
\]
Then
\[
 \ddc u_1\wedge\cdots\wedge\ddc u_n
                  \geq\left(\prod_{i=1}^n f_i^{1/n}\right)\eta^n.
\]
\end{corollary}
\begin{proof}
For $n=1$ this is the hypothesis, so let $n\geq2$. Put $u_{i,\varepsilon}=u_i*\rho_\varepsilon$, with the same radial mollifiers for all $i$. By the last assertion of \cref{prop:bounded-monge-ampere-linear-density}, applied to each $u_i$ with the density $f_i$,
\[
 \det\mathcal H(u_{i,\varepsilon})\geq(f_i^{1/n}*\rho_\varepsilon)^n\qquad(1\leq i\leq n).
\]
By \eqref{eq:bt-determinant-normalization}, applied to the smooth function $\sum_it_iu_{i,\varepsilon}$, the exterior powers of smooth forms satisfy
\[
 \Bigl(\ddc\sum_{i=1}^nt_iu_{i,\varepsilon}\Bigr)^n=\det\Bigl(\sum_{i=1}^nt_i\mathcal H(u_{i,\varepsilon})\Bigr)\eta^n\qquad(t_1,\ldots,t_n\in\mathbb R).
\]
Both sides are polynomials in $t_1,\ldots,t_n$. Since $2$-forms commute, the coefficient of $t_1\cdots t_n$ on the left is $n!\,\ddc u_{1,\varepsilon}\wedge\cdots\wedge\ddc u_{n,\varepsilon}$, and on the right it is $n!\,D(\mathcal H(u_{1,\varepsilon}),\ldots,\mathcal H(u_{n,\varepsilon}))\eta^n$. Since the Bedford--Taylor product of smooth functions is the exterior product of the forms $\ddc u_{i,\varepsilon}$ (\cref{prop:bt-positive-recursion}(1)), we obtain
\[
 \ddc u_{1,\varepsilon}\wedge\cdots\wedge\ddc u_{n,\varepsilon}
 =D\bigl(\mathcal H(u_{1,\varepsilon}),\ldots,\mathcal H(u_{n,\varepsilon})\bigr)\eta^n
 \geq\prod_{i=1}^n(f_i^{1/n}*\rho_\varepsilon)\,\eta^n,
\]
where the inequality is \cref{lem:multifactor-mixed-discriminant}, applied to the complex Hessians of the smooth plurisubharmonic functions $u_{i,\varepsilon}$, which are positive semidefinite (\cref{prop:smooth-psh-criterion}). As $\varepsilon\downarrow0$, the density on the right converges locally uniformly to $\prod_if_i^{1/n}$, and the left side converges weakly to $\ddc u_1\wedge\cdots\wedge\ddc u_n$ by \cref{thm:bt-decreasing-continuity}, since $u_{i,\varepsilon}\downarrow u_i$ for each $i$ (\cref{prop:smooth-psh-criterion}).
\end{proof}
\index{Bedford--Taylor product!multifactor density inequality}


\section{The ball Dirichlet problem with continuous density}
\label{sec:bt-continuous-ball-dirichlet}

The linear density inequalities make subsolutions stable under maxima and controlled quadratic perturbations. These operations allow the homogeneous ball envelope to be replaced by an envelope with prescribed density. The exact-center automorphisms still give bounded second differences, provided their Jacobian factor is included in the density.

\begin{lemma}\label{lem:continuous-density-subsolution-algebra}
Let $u,v\in\PSH(\Omega)\cap L^\infty_{\loc}$ and let $F,G\ge0$ be continuous with $(\ddc u)^n\ge F^n\eta^n$ and $(\ddc v)^n\ge G^n\eta^n$. For $0\le t\le1$ and $a\ge0$,
\[
 \bigl(\ddc((1-t)u+tv+a|z|^2)\bigr)^n
 \ge\bigl((1-t)F+tG+a\bigr)^n\eta^n.
\]
Also $(\ddc\max\{u,v\})^n\ge\min\{F,G\}^n\eta^n$.
\end{lemma}
\begin{proof}
For every admissible $B$, the linear inequalities in \cref{prop:bounded-monge-ampere-linear-density} give
\[
 L_B((1-t)u+tv+a|z|^2)
 \ge n((1-t)F+tG)+a\operatorname{tr}B
 \ge n((1-t)F+tG+a).
\]
Here $\operatorname{tr}B\ge n$ by arithmetic--geometric mean. The same proposition gives the first assertion. An upper test for $\max\{u,v\}$ at a point touches at least one of $u,v$ there. Its determinant is at least the corresponding density, and hence at least $\min\{F,G\}^n$. Apply the upper-test characterization again.
\end{proof}

\begin{theorem}\label{thm:continuous-density-ball-dirichlet}
Let $B=B(a,R)\subset\mathbb C^n$, $n\ge1$. For $f\in C(\overline B)$ with $f\ge0$ and real $\varphi\in C(\partial B)$, there is a unique $H_{f,\varphi}\in\PSH(B)\cap C(\overline B)$ with
\[
 (\ddc H_{f,\varphi})^n=f\eta^n,
 \qquad H_{f,\varphi}|_{\partial B}=\varphi.
\]
It is the pointwise supremum of bounded plurisubharmonic $v$ satisfying $(\ddc v)^n\ge f\eta^n$ and upper boundary limits at most $\varphi$. For any other such data $g,\chi$,
\begin{equation}\label{eq:ball-dirichlet-root-stability}
 \|H_{f,\varphi}-H_{g,\chi}\|_{L^\infty(B)}
 \le\|\varphi-\chi\|_{L^\infty(\partial B)}
       +R^2\|f^{1/n}-g^{1/n}\|_{L^\infty(B)}.
\end{equation}
If $f^{1/n}$ and $\varphi$ have real $C^2$ extensions to neighborhoods of $\overline B$ and $\partial B$, respectively, then $H_{f,\varphi}\in C^{1,1}_{\loc}(B)$.
\end{theorem}
\begin{proof}
An affine change $z=a+R\zeta$ reduces the problem to the unit ball. The transformed density is $R^{2n}f(a+R\zeta)$, so its root is $R^2f^{1/n}(a+R\zeta)$. We prove the assertions there and put $\psi=|z|^2-1$ and $F=f^{1/n}$.

\emph{Step 1. The subsolution envelope for regular data.}
First suppose $F$ and $\varphi$ have the stated $C^2$ extensions. Extend $\varphi$ as a real $C^2$ function $b$ on a neighborhood of $\overline B$. For a sufficiently large $A$, $s=b+A\psi$ is strictly plurisubharmonic and $\det H(s)\ge f$ on $B$. It is a bounded subsolution with the required boundary value. Every bounded competitor is at most the harmonic extension $P\varphi$, by \cref{thm:subharmonic-comparison}. Let $H$ be their supremum. By \cref{lem:continuous-density-subsolution-algebra}, finite maxima of competitors, including $s$, remain competitors. Choquet selection in \cref{thm:psh-upper-envelopes} supplies such finite maxima increasing almost everywhere to the upper regularization $H^*$. They lie between $s$ and $P\varphi$. Increasing continuity in \cref{thm:bt-increasing-continuity} therefore gives $(\ddc H^*)^n\ge f\eta^n$. The same two barriers give boundary value $\varphi$, so $H^*$ is a competitor. Hence $H=H^*$ pointwise, and $H$ is bounded and plurisubharmonic.

\emph{Step 2. Second differences and equality of the density.}
Fix a compact set of centers $K\Subset B$ and use $S_{a,h}$ from \cref{lem:ball-centered-automorphisms}. Write $J_{a,h}(z)=\det_{\mathbb C}D_zS_{a,h}$. Local smooth approximation and decreasing continuity show that Bedford--Taylor measures commute with these biholomorphisms. Consequently
\[
 \bigl(\ddc(H\circ S_{a,h})\bigr)^n
 \ge (F\circ S_{a,h})^n|J_{a,h}|^2\eta^n.
\]
The explicit formula for the automorphisms has denominators bounded away from zero on $\overline B$, uniformly for $a\in K$ and sufficiently small $h$. Its inverse has the same property. Thus $|J_{a,h}|^{2/n}$ is smooth and strictly positive there, with its first two real $h$ derivatives uniformly bounded. Taylor's formula and $S_{a,0}=\mathrm{id}$ give constants $C_1,C_2\ge0$ such that
\[
 \begin{aligned}
 \tfrac12\bigl((F\circ S_{a,h})|J_{a,h}|^{2/n}
 +(F\circ S_{a,-h})|J_{a,-h}|^{2/n}\bigr)
 &\ge F-C_1|h|^2 &&\text{on }\overline B,\\
 \tfrac12(\varphi\circ S_{a,h}+\varphi\circ S_{a,-h})
 &\le\varphi+C_2|h|^2 &&\text{on }\partial B.
 \end{aligned}
\]
The constants depend on $K,n$ and the fixed $C^2$ data, independently of $a,h,z$. By \cref{lem:continuous-density-subsolution-algebra},
\[
 V_{a,h}=\tfrac12(H\circ S_{a,h}+H\circ S_{a,-h})
            +C_1|h|^2\psi-C_2|h|^2
\]
has Monge--Amp\`ere measure at least $f\eta^n$ and upper boundary limits at most $\varphi$. Indeed the average's root lower bound is the nonnegative average in the first line above; adding $C_1|h|^2\psi$ raises it by $C_1|h|^2$. Therefore $V_{a,h}\le H$. Evaluating at $a$ and using $-1\le\psi(a)\le0$ gives
\[
 H(a+h)+H(a-h)-2H(a)\le2(C_1+C_2)|h|^2.
\]
As in the proof of \cref{thm:homogeneous-ball-dirichlet}, convolution turns this into an upper bound for every real directional Hessian; plurisubharmonicity in the direction $\mathrm ih$ gives its lower bound. \Cref{prop:bounded-hessian-taylor} then gives the canonical locally $C^{1,1}$ representative and its almost-everywhere second-order Taylor expansion. Bounded Hessian coefficients, Lebesgue differentiation and uniform Bedford--Taylor continuity identify $(\ddc H)^n=\det(H_{j\bar k})\eta^n$.

At a Taylor point $x$ where $\det(H_{j\bar k})(x)>f(x)$, the complex Hessian is positive definite. Let $Q$ be the real quadratic Taylor polynomial. Choose $\varepsilon>0$ so small that $Q-\varepsilon|z-x|^2$ remains strictly plurisubharmonic and its determinant exceeds $f$ on a small ball about $x$. Shrink its radius $r$ so that the Taylor remainder on the sphere is less than $\varepsilon r^2/2$. Then
$q=Q-\varepsilon|z-x|^2+\varepsilon r^2/4$ is below $H$ on that sphere and above it at $x$. Paste $\max\{H,q\}$ inside the small ball and $H$ outside. Gluing and \cref{lem:continuous-density-subsolution-algebra} give a bounded competitor exceeding $H$, a contradiction. Since $\det(H_{j\bar k})\ge f$ almost everywhere already, equality holds almost everywhere. This proves the equation for regular data.

\emph{Step 3. Root-density stability and continuous data.}
For two regular data sets, put $\delta=\|\varphi-\chi\|_\infty$ and $\varepsilon=\|f^{1/n}-g^{1/n}\|_\infty$. The function $H_{g,\chi}+\varepsilon\psi-\delta$ has upper boundary limits at most $\varphi$ and, by \cref{lem:continuous-density-subsolution-algebra}, measure at least $(g^{1/n}+\varepsilon)^n\eta^n\ge f\eta^n$. It is a competitor for $H_{f,\varphi}$. Reverse the data sets to obtain \eqref{eq:ball-dirichlet-root-stability} with $R=1$.

For continuous data, extend $F\ge0$ continuously outside $\overline B$ by radial projection and convolve with nonnegative kernels. This gives smooth nonnegative $F_j\to F$ uniformly on $\overline B$. Smooth boundary approximations $\varphi_j\to\varphi$ are obtained on an annular neighborhood. The estimate just proved makes $H_{F_j^n,\varphi_j}$ uniformly Cauchy on $\overline B$. Its limit $U$ is continuous and plurisubharmonic, with boundary value $\varphi$. Uniform continuity of the bounded products gives $(\ddc U)^n=f\eta^n$. Applying the estimate to approximations of two sets of data and taking their uniform limits proves the stated stability bound for their solutions.

\emph{Step 4. Uniqueness and the envelope.}
Let $U$ be the solution constructed above and let $v$ be any bounded plurisubharmonic competitor for these continuous data. For $\varepsilon>0$, put $v_\varepsilon=v+\varepsilon\psi$. Its measure is at least $(F+\varepsilon)^n\eta^n$, whereas that of $U$ is $F^n\eta^n$. Boundary comparison gives
\[
 \int_{\{U<v_\varepsilon\}}(F+\varepsilon)^n\eta^n
 \le\int_{\{U<v_\varepsilon\}}F^n\eta^n.
\]
The strict difference makes this Borel set Lebesgue null. Uniqueness of canonical plurisubharmonic representatives gives $U\ge v_\varepsilon$ pointwise, as in \cref{thm:bt-comparison}. Let $\varepsilon\downarrow0$ to get $U\ge v$. Since $U$ itself is a competitor, it equals their pointwise envelope. Any other continuous solution is a competitor and is therefore at most $U$; applying the same argument with the solutions interchanged gives equality. Scaling back proves all assertions.
\end{proof}
\index{Dirichlet problem!continuous Monge--Amp\`ere density}


\section{Bounded measurable densities}
\label{sec:bt-measurable-ball-dirichlet}

Uniform convergence of densities is too restrictive for approximation by local averages. A common bound on their densities instead controls one side of the difference between solutions by capacity. Hartogs' lemma and quasi-continuity control the other side, so bounded-product continuity still identifies the limit.

\begin{lemma}\label{lem:local-comparison-capacity-sublevels}
Let $D\subset\mathbb C^n$ be a bounded domain and let $u,v$ be bounded plurisubharmonic functions on $D$. If $s,t>0$ and $\{u<v-s\}\Subset D$, then
\[
 t^n\operatorname{Cap}_D(\{u<v-s-t\})
 \le\int_{\{u<v-s\}}(\ddc u)^n.
\]
Here $E\Subset D$ means that its closure is a compact subset of $D$.
\end{lemma}
\begin{proof}
For $w\in\PSH(D)$ with $-1\le w\le0$,
\[
 \{u<v-s-t\}\subset\{u<v-s+tw\}\subset\{u<v-s\}.
\]
The middle comparison set is relatively compact. Positivity and expansion of the bounded products give
$\bigl(\ddc(v-s+tw)\bigr)^n\ge t^n(\ddc w)^n$.
Apply \cref{thm:bt-comparison} to $u$ and $v-s+tw$, and take the supremum over $w$. Translating these tests by $1$ gives exactly the capacity tests in \cref{def:monge-ampere-capacity}.
\end{proof}

\begin{theorem}\label{thm:bounded-density-ball-dirichlet}
Let $B=B(a,R)\subset\mathbb C^n$, $n\ge1$, let $0\le f\in L^\infty(B)$, and let $\varphi\in C(\partial B)$ be real. There is a unique bounded $U_{f,\varphi}\in\PSH(B)$ such that
\[
 (\ddc U_{f,\varphi})^n=f\eta^n,
 \qquad\lim_{B\ni z\to x}U_{f,\varphi}(z)=\varphi(x)
       \quad(x\in\partial B).
\]
It is the pointwise envelope of bounded plurisubharmonic subsolutions with upper boundary limits at most $\varphi$. The stability estimate \eqref{eq:ball-dirichlet-root-stability} holds also for these bounded measurable densities.

If $0\le f_j\le M$ almost everywhere for a fixed finite $M$, $f_j\to f$ in $L^1(B)$ and $\varphi_j\to\varphi$ uniformly on $\partial B$, then
\[
 U_{f_j,\varphi_j}\longrightarrow U_{f,\varphi}
 \quad\hbox{in }L^1(B)\hbox{ and locally in capacity}.
\]
\end{theorem}
\begin{proof}
\emph{Step 1. Bounded approximations and their boundary barriers.}
Extend $f$ by zero outside $B$ and convolve with nonnegative kernels. The resulting smooth $f_j$ satisfy $0\le f_j\le M=\|f\|_\infty$ and $f_j\to f$ in $L^1(B)$, by \cref{lem:mollifier-density}. Solve the continuous-density problems by \cref{thm:continuous-density-ball-dirichlet}, with the fixed boundary value $\varphi$, and call their solutions $u_j$. Comparison gives
\[
 H_{M,\varphi}\le u_j\le H_{0,\varphi}.
\]
For example, to justify either ordering, add $\varepsilon(|z-a|^2-R^2)$ to the proposed smaller function. Positivity of mixed products raises its measure by at least $\varepsilon^n\eta^n$. Comparison makes the strict reversed-order set Lebesgue null; canonical representatives give the pointwise ordering as in \cref{thm:bt-comparison}, and then $\varepsilon\downarrow0$. Both displayed barriers are continuous on $\overline B$ and equal $\varphi$ on its boundary.

By \cref{thm:psh-compactness-hartogs}, a subsequence converges in $L^1_{\loc}(B)$ to a bounded plurisubharmonic $U$. The uniform bounds exclude collapse to $-\infty$. Canonical representatives preserve the barrier inequalities, so $U$ has the required boundary limits. Uniform boundedness also promotes convergence to $L^1(B)$: first discard a boundary collar of small volume, then use the local convergence on its complement.

\emph{Step 2. Capacity convergence.}
Fix $K\Subset B$ and $\delta,\varepsilon>0$. Quasi-continuity gives an open $G\subset B$ with $\operatorname{Cap}_B(G)<\varepsilon$ such that $U$ is continuous on $B\setminus G$. On the compact $K\setminus G$, Hartogs' lemma gives $u_j-U<\delta$ for all sufficiently large $j$. Indeed upper semicontinuity bounds $U$ above near each point by its value there plus $\delta/4$; continuity on $K\setminus G$ bounds its values below there by that value minus $\delta/4$. Apply Hartogs' lemma with a constant majorant on a smaller closed ball and use a finite cover. Thus
\[
 \limsup_j\operatorname{Cap}_B(K\cap\{u_j>U+\delta\})
 \le\varepsilon.
\]
The upper-difference set on $K$ is contained in the open set $G$ for all sufficiently large $j$, so the same bound holds for its outer capacity. Let $\varepsilon\downarrow0$.

The two continuous boundary barriers show that $\{u_j<U-\delta\}$ lies in a fixed relatively compact subset of $B$, independently of $j$. With $s=t=\delta$, \cref{lem:local-comparison-capacity-sublevels} gives
\[
 \begin{split}
 \delta^n\operatorname{Cap}_B(\{u_j<U-2\delta\})
 &\le\int_{\{u_j<U-\delta\}}f_j\eta^n\\
 &\le \frac M\delta\int_B(U-u_j)_+\eta^n\longrightarrow0.
 \end{split}
\]
Here $\eta^n=n!\pi^{-n}\,\mathrm dV$. The bounded difference $U-u_j$ is quasi-continuous. Applying \cref{lem:quasi-continuous-level-capacity} on a compact neighborhood of $K$ converts this last estimate into vanishing outer capacity on $K$. Together with the upper-difference estimate, this proves local convergence in the ordered sense of \cref{def:local-capacity-convergence}.

\emph{Step 3. The equation, envelope and uniqueness.}
Capacity continuity in \cref{thm:bt-capacity-continuity} gives $(\ddc U)^n=f\eta^n$. Let $v$ be any bounded subsolution with upper boundary limits at most $\varphi$, and put $v_\varepsilon=v+\varepsilon(|z-a|^2-R^2)$. Positivity of mixed products gives
\[
 (\ddc v_\varepsilon)^n\ge f\eta^n+\varepsilon^n\eta^n.
\]
Comparison with $U$ makes $\{U<v_\varepsilon\}$ Lebesgue null. The integrals of $f\eta^n$ are finite, so their cancellation is legitimate. Canonical representatives give $U\ge v_\varepsilon$ pointwise. Letting $\varepsilon\downarrow0$ proves $U\ge v$. The solution itself is a subsolution, which identifies the pointwise envelope. Interchanging two bounded solutions proves uniqueness.

\emph{Step 4. Stability under a common density bound.}
For any sequence of bounded measurable $f_j$ with the stated common bound and fixed boundary value $\varphi$, the same barriers give $L^1$ subsequential compactness. Steps 2 and 3 apply unchanged, since their only density estimate was $(\ddc u_j)^n\le M\eta^n$. Every subsequential limit solves the equation for $f$, and hence equals $U_{f,\varphi}$. Every subsequence has such a further subsequence, proving convergence of the whole sequence in $L^1(B)$ and locally in capacity.

For fixed density, comparison and addition of constants give
\[
 \|U_{f,\varphi}-U_{f,\chi}\|_\infty
 \le\|\varphi-\chi\|_\infty.
\]
Apply this with $f=f_j$ to reduce varying boundary data to the fixed-boundary assertion. Uniformly vanishing differences also vanish locally in capacity, since compact subsets have finite capacity.

Finally take two bounded measurable densities $f,g$. Extend their nonnegative roots by zero and convolve with the same kernels. The smooth roots $F_j,G_j$ satisfy
$\|F_j-G_j\|_\infty\le\|f^{1/n}-g^{1/n}\|_\infty$, and their powers converge to $f,g$ in $L^1(B)$ with a common bound. The continuous-density stability estimate applies to the two approximating solutions. Step 4 gives $L^1$ limits; choose a common almost-everywhere convergent subsequence and pass the estimate to the limit. Each resulting inequality holds pointwise by uniqueness of canonical plurisubharmonic representatives after addition of its constant. This proves \eqref{eq:ball-dirichlet-root-stability} in the measurable case as well.
\end{proof}
\index{Dirichlet problem!bounded measurable density}


\section{Local energy moments}
\label{sec:local-energy-moments}
In approximation arguments a common bound on the densities may be unavailable. A bounded subsolution instead controls the measure's integrals against zero-boundary potentials. The truncation argument below derives that control from B\l{}ocki's difference estimate, without introducing unbounded Monge--Amp\`ere products. Sources are discussed in \cref{sec:history-bounded-products}.

\begin{definition}\label{def:bounded-zero-boundary-energy}
Let $B\subset\mathbb C^n$, $n\geq1$, be a bounded ball. The class $\mathcal E_0(B)$ consists of bounded plurisubharmonic functions $\varphi\leq0$ such that
\[
 \lim_{z\to\zeta,\ z\in B}\varphi(z)=0\quad(\zeta\in\partial B),
 \qquad \int_B(\ddc\varphi)^n<\infty.
\]
We include the zero function. For $p\geq1$ its local energy moment is
\[
 e_p(\varphi)=\int_B(-\varphi)^p(\ddc\varphi)^n.
\]
\end{definition}
Boundedness and finite mass make every such moment finite. Compactness of $\partial B$ makes the displayed boundary limit uniform near the boundary. In particular each set $\{\varphi\leq-s\}$, $s>0$, lies in a compact subset of $B$.
\index{local energy!bounded zero-boundary class}

\begin{theorem}\label{thm:bounded-subsolution-local-energy-moments}
Let $\mu$ be a finite positive measure on $B$ with
$\mu\leq(\ddc w)^n$ for some bounded $w\in\PSH(B)$. Put
\[
 m=\mu(B),\qquad M=\sup_Bw-\inf_Bw,
 \qquad K_{n,p}=2^{n+p}(n+1)!M^n.
\]
For every $\varphi\in\mathcal E_0(B)$, $p\geq1$ and $t>0$,
\begin{equation}\label{eq:local-energy-measure-tail}
 \mu(\{\varphi<-t\})
         \leq\min\{m,K_{n,p}e_p(\varphi)t^{-(n+p)}\}.
\end{equation}
Consequently,
\begin{equation}\label{eq:local-energy-measure-moment}
 \int_B(-\varphi)^p\,\mathrm d\mu
 \leq\left(1+\frac pn\right)
      K_{n,p}^{p/(n+p)}m^{n/(n+p)}
                          e_p(\varphi)^{p/(n+p)}.
\end{equation}
In particular $\int_B(-\varphi)\,\mathrm d\mu\leq A e_1(\varphi)^{1/(n+1)}$, where $A$ depends only on $n,M,m$.
\end{theorem}
\begin{proof}
For $s>0$ let $h_s=\max\{\varphi,-s\}$ and
$d_s=h_s-\varphi=(-\varphi-s)_+$. Choose a smaller concentric ball $B_s\Subset B$ containing $\{\varphi\leq-s\}$ in its interior and with $\varphi>-s$ on a collar of its boundary. Then $h_s=\varphi$ on that collar. All three functions $\varphi,h_s,w$ are bounded plurisubharmonic near $\overline{B_s}$. Applying \cref{prop:blocki-difference-estimate} there, and using its oscillation bound $\operatorname{osc}_{B_s}w\leq M$, gives
\[
 \int_B d_s^{n+1}\,\mathrm d\mu
  \leq\int_{B_s}d_s^{n+1}(\ddc w)^n
  \leq(n+1)!M^n\int_{B_s}d_s(\ddc\varphi)^n.
\]
The first integral is over $B$ because $d_s$ vanishes outside $B_s$. Set $s=t/2$. On $\{\varphi<-t\}$ we have $d_s>t/2$, while
\[
 \int_{B_s}d_s(\ddc\varphi)^n
  \leq\int_{\{-\varphi>t/2\}}(-\varphi)(\ddc\varphi)^n
  \leq(t/2)^{1-p}e_p(\varphi).
\]
The last inequality uses $p\geq1$. Combining the two estimates proves the second tail bound; the first is the total-mass bound.

If $m=0$ the moment estimate is immediate. If $K_{n,p}e_p(\varphi)=0$, the tail bound makes every positive-depth sublevel set $\mu$-null, so the moment is again zero. Otherwise put
\[
 T=\left(\frac{K_{n,p}e_p(\varphi)}m\right)^{1/(n+p)}.
\]
Layer integration and \eqref{eq:local-energy-measure-tail} give
\[
 \begin{aligned}
 \int_B(-\varphi)^p\,\mathrm d\mu
 &=p\int_0^\infty t^{p-1}\mu(\{\varphi<-t\})\,\mathrm dt\\
 &\leq mT^p
       +pK_{n,p}e_p(\varphi)\int_T^\infty t^{-n-1}\,\mathrm dt\\
 &=mT^p+\frac pnK_{n,p}e_p(\varphi)T^{-n}.
 \end{aligned}
\]
Substitution of $T$ proves \eqref{eq:local-energy-measure-moment}. Taking $p=1$ gives the final assertion. The argument applies also when $n=1$ and requires no boundary regularity of $w$.
\end{proof}
\index{local energy!measure moment estimate}
\index{Monge--Amp\`ere measure!bounded subsolution}

\Needspace{20\baselineskip}
\begin{proposition}\label{prop:canonical-grid-energy-control}
Suppose in the preceding theorem that $\mu$ has compact support $K\Subset B$. Extend $\mu$ by zero outside $B$. Partition $\mathbb R^{2n}$ into half-open cubes $Q$ of side length $h_j\downarrow0$, and, for sufficiently small $h_j$, set
\[
 \mu_j=\sum_Q\frac{\mu(Q)}{|Q|}\mathbf1_Q\,\mathrm dV.
\]
Then $\mu_j(B)=m$, each $\mu_j$ has bounded density and compact support in $B$, and $\mu_j\to\mu$ weakly. For $p\geq1$ there is $D_{n,p}>0$, depending only on $n,p$, such that
\begin{equation}\label{eq:canonical-grid-energy-moment}
 \int_B(-\varphi)^p\,\mathrm d\mu_j
        \leq D_{n,p}A_p e_p(\varphi)^{p/(n+p)},
                   \qquad\varphi\in\mathcal E_0(B),
\end{equation}
where $A_p=(1+p/n)K_{n,p}^{p/(n+p)}m^{n/(n+p)}$.
The zero-boundary solutions $U_j$ of $(\ddc U_j)^n=\mu_j$ satisfy
\[
 e_p(U_j)\leq(D_{n,p}A_p)^{(n+p)/n},\qquad p\geq1.
\]
They are relatively compact in $L^1(B)$, and each $L^1$ limit has a canonical nonpositive plurisubharmonic representative.
\end{proposition}
\begin{proof}
If $m=0$, take $\mu_j=0$ and $U_j=0$; all assertions follow. Otherwise put $d=\operatorname{dist}(K,\partial B)>0$ and require $h_j<d/(6\sqrt{2n})$. Every cube with positive $\mu$-mass meets $K$ and lies in $B$. Only finitely many such cubes occur. Their masses sum to $m$, and their bounded densities are constant on the cubes. For a continuous test function, averaging over the cube containing a point changes its value by at most its modulus of continuity at $\sqrt{2n}h_j$, on a fixed compact neighborhood of $K$. Integration against $\mu$ proves weak convergence.

Fix a cube $Q$ with $\mu(Q)>0$ and $z\in Q$. The polydisc centered at $z$ with coordinate radius $4\sqrt2h_j$ lies in $B$, and its half-sized polydisc contains $Q$. Translation and scaling in \cref{prop:negative-psh-polydisc-lp} give
\[
 \frac1{|Q|}\int_Q(-\varphi)^p\,\mathrm dV
                         \leq D_{n,p}(-\varphi(z))^p.
\]
The constant is independent of $Q,j,z$; the ratio between the polydisc volume scale and $|Q|=h_j^{2n}$ is fixed. Integrate this inequality over $z\in Q$ against $\mu$, then sum over the cubes. This gives
\[
 \int_B(-\varphi)^p\,\mathrm d\mu_j
                  \leq D_{n,p}\int_B(-\varphi)^p\,\mathrm d\mu.
\]
The preceding theorem proves \eqref{eq:canonical-grid-energy-moment}.

Each bounded-density solution exists by \cref{thm:bounded-density-ball-dirichlet}. Its zero boundary limit and mass $m$ put it in $\mathcal E_0(B)$. Substitute $\varphi=U_j$ in \eqref{eq:canonical-grid-energy-moment}. If its energy is positive, divide by $e_p(U_j)^{p/(n+p)}$; if it is zero, the stated bound is immediate.

Apply the preceding theorem once more with the fixed measure $(\ddc\rho)^n$, where $\rho(z)=|z-a|^2-R^2$ on $B=B(a,R)$. This measure is a fixed positive multiple of Euclidean volume. The uniform $e_2(U_j)$ bound therefore gives a uniform $L^2(B)$ bound. Local psh compactness, \cref{thm:psh-compactness-hartogs}, yields a subsequence converging in $L^1$ on compact subsets to a nonpositive psh limit; the $L^2$ bound excludes collapse to $-\infty$. Pass to a further subsequence converging almost everywhere on $B$. Fatou gives the same $L^2$ bound for the limit. On a boundary strip, Cauchy--Schwarz bounds the $L^1$ error by a uniform constant times the square root of its volume. Exhausting $B$ by compact subsets now gives convergence in $L^1(B)$. The same argument applies to every subsequence.
\end{proof}
\index{Monge--Amp\`ere measure!canonical grid approximation}


\subsection{Convergence of energy integrals}
\label{sec:canonical-grid-energy-convergence}
We shall identify the limit $u$ of a convergent subsequence of canonical approximants as a solution of $(\ddc u)^n=\mu$. For this we must pass to the limit in integrals such as $\int U_j\,\mathrm d\mu_j$, in which both the integrand and the measure depend on $j$, and the energy bounds of \cref{prop:canonical-grid-energy-control} control only the large negative values of the integrand. We first prove that $L^1_{\loc}$ convergence of uniformly bounded plurisubharmonic functions implies $L^1$ convergence against any fixed finite measure that does not charge sets of outer capacity zero (\cref{lem:bounded-psh-fixed-measure-continuity}). Writing integrals against $\mu_j$ as integrals of cube averages against $\mu$, we then transfer this continuity to the varying measures $\mu_j$ (\cref{thm:canonical-grid-self-integral-convergence}).

\begin{lemma}\label{lem:bounded-psh-fixed-measure-continuity}
Let $\nu$ be a finite positive Borel measure on a ball $B$ such that $\nu(E)=0$ for every Borel set $E\subset B$ satisfying
\[
 \operatorname{Cap}_B^*(E\cap K)=0\qquad\text{for every compact }K\subset B.
\]
Let $L\geq0$, let $v_j$ be plurisubharmonic functions on $B$ with $-L\leq v_j\leq0$, and suppose that $v_j\to v$ in $L^1_{\loc}(B)$, where $v$ is the plurisubharmonic representative of the limit. Then
\[
 \int_B|v_j-v|\,\mathrm d\nu\longrightarrow0.
\]
The hypothesis on $\nu$ holds if $\nu\leq(\ddc w)^n$ for a bounded plurisubharmonic function $w$ on $B$.
\end{lemma}
\begin{proof}
The bounds $-L\le v\le0$ hold almost everywhere, hence everywhere by the uniqueness assertion of \cref{thm:psh-distribution-representative}.

\emph{Step 1. The upper limit of $v_j$.} We show that
\begin{equation}\label{eq:fixed-measure-psh-limsup}
 \limsup_{j\to\infty}v_j=v\qquad\nu\text{-almost everywhere}.
\end{equation}
For $\ell\ge1$ let
\[
 b_\ell=\left(\sup_{j\geq\ell}v_j\right)^*
\]
be the upper semicontinuous regularization of the supremum. By \cref{thm:psh-upper-envelopes}, these are plurisubharmonic functions with values in $[-L,0]$, and they decrease as $\ell$ increases. We claim that $b_\ell\downarrow v$ everywhere. Let $z\in B$ and $\varepsilon>0$. By upper semicontinuity, $v<v(z)+\varepsilon$ on a neighborhood of a closed ball $\overline{B'}$ centered at $z$, and Hartogs' lemma (\cref{thm:psh-compactness-hartogs}), applied with the constant majorant $v(z)+\varepsilon$, gives $v_j<v(z)+\varepsilon$ on $\overline{B'}$ for all large $j$. Hence $b_\ell(z)\le v(z)+\varepsilon$ for all large $\ell$, and the decreasing limit of $b_\ell$ is at most $v$. Conversely, $b_\ell\ge v_j$ for all $j\ge\ell$, and a subsequence of $(v_j)$ converges to $v$ almost everywhere, so $b_\ell\ge v$ almost everywhere. The decreasing limit of the $b_\ell$ is plurisubharmonic by \cref{thm:psh-decreasing-nets} and equals $v$ almost everywhere, hence everywhere by uniqueness of plurisubharmonic representatives.

The sets $N_\ell=\{\sup_{j\ge\ell}v_j<b_\ell\}$ are Borel, and by \cref{cor:psh-envelope-zero-capacity} they satisfy $\operatorname{Cap}_B^*(N_\ell\cap K)=0$ for every compact $K\subset B$. By countable subadditivity of outer capacity, the union $N=\bigcup_\ell N_\ell$ has the same property, so $\nu(N)=0$ by hypothesis. For $z\notin N$ we have $\sup_{j\ge\ell}v_j(z)=b_\ell(z)$ for every $\ell$, and therefore $\limsup_jv_j(z)=\lim_\ell b_\ell(z)=v(z)$. This proves \eqref{eq:fixed-measure-psh-limsup}. The argument applies to every sequence of plurisubharmonic functions with values in $[-L,0]$ that converges to $v$ in $L^1_{\loc}(B)$.

\emph{Step 2. Convergence of $\int v_j\,\mathrm d\nu$.} Let $\mathcal C_\ell$ be the set of finite convex combinations of $v_\ell,v_{\ell+1},\ldots$, and let
\[
 d_\ell=\inf_{r\in\mathcal C_\ell}\|r\|_{L^2(\nu)}.
\]
Since $\mathcal C_{\ell+1}\subset\mathcal C_\ell$, the numbers $d_\ell$ increase, and they are bounded by $L\,\nu(B)^{1/2}$. Choose $r_\ell\in\mathcal C_\ell$ with
\[
 \|r_\ell\|_{L^2(\nu)}^2\leq d_\ell^2+\ell^{-1}.
\]
For $k\geq\ell$ the midpoint $(r_k+r_\ell)/2$ belongs to the convex set $\mathcal C_\ell$, so its norm is at least $d_\ell$, and the parallelogram identity gives
\[
 \|r_k-r_\ell\|_{L^2(\nu)}^2
 =2\|r_k\|_{L^2(\nu)}^2+2\|r_\ell\|_{L^2(\nu)}^2-4\Bigl\|\frac{r_k+r_\ell}2\Bigr\|_{L^2(\nu)}^2
 \leq2(d_k^2-d_\ell^2)+2k^{-1}+2\ell^{-1}.
\]
Hence $(r_\ell)$ is a Cauchy sequence in $L^2(\nu)$, and it converges in $L^2(\nu)$ to a function $h$ by completeness of $L^2(\nu)$. (A subsequence with summable successive distances converges $\nu$-almost everywhere and in $L^2(\nu)$, by the triangle inequality and monotone convergence applied to the sum of the absolute differences; the Cauchy property then gives convergence of the whole sequence.) Some subsequence of $(r_\ell)$ converges to $h$ $\nu$-almost everywhere. On the other hand, the $r_\ell$ are plurisubharmonic with values in $[-L,0]$, and $r_\ell\to v$ in $L^1_{\loc}(B)$, since $\|r_\ell-v\|_{L^1(K)}\le\sup_{j\ge\ell}\|v_j-v\|_{L^1(K)}$ for every compact $K\subset B$. Applying \eqref{eq:fixed-measure-psh-limsup} to that subsequence, we find $h=v$ $\nu$-almost everywhere. Thus $r_\ell\to v$ in $L^2(\nu)$, and hence $\int r_\ell\,\mathrm d\nu\to\int v\,\mathrm d\nu$.

We now show that $\int v_j\,\mathrm d\nu\to\int v\,\mathrm d\nu$. Otherwise there are $c>0$ and a subsequence $(v_{j_i})$ such that either $\int v_{j_i}\,\mathrm d\nu\ge\int v\,\mathrm d\nu+c$ for all $i$, or $\int v_{j_i}\,\mathrm d\nu\le\int v\,\mathrm d\nu-c$ for all $i$. The subsequence satisfies the hypotheses of the lemma, so the preceding paragraph provides convex combinations of its tails whose integrals converge to $\int v\,\mathrm d\nu$. This is impossible, since every convex combination of the $v_{j_i}$ satisfies the same inequality as the $v_{j_i}$ themselves.

\emph{Step 3. Conclusion.} By Hartogs' lemma, $\limsup_jv_j\le v$ everywhere, so $(v_j-v)_+\to0$ pointwise. These functions are bounded by $L$, so $\int(v_j-v)_+\,\mathrm d\nu\to0$ by dominated convergence. Since
\[
 |v_j-v|=2(v_j-v)_+-(v_j-v),
\]
Step~2 proves the assertion.

Finally, suppose that $\nu\le(\ddc w)^n$ with $w$ bounded and plurisubharmonic, and put $M=\sup_Bw-\inf_Bw$. If $M=0$, then $w$ is constant and $\nu=0$. Otherwise $(w-\inf_Bw)/M$ belongs to the family defining $\operatorname{Cap}_B$, so $\nu(A)\leq M^n\operatorname{Cap}_B(A)\le M^n\operatorname{Cap}_B^*(A)$ for every Borel set $A\subset B$. If $E$ is as in the hypothesis, then $\nu(E\cap K)=0$ for every compact $K\subset B$, and exhausting $B$ by compact sets gives $\nu(E)=0$.
\end{proof}
\index{Monge--Amp\`ere measure!fixed-measure continuity}

\begin{theorem}\label{thm:canonical-grid-self-integral-convergence}
Let $\mu$, $\mu_j$ and $U_j$ be as in \cref{prop:canonical-grid-energy-control}. Suppose that a subsequence of $(U_j)$, relabeled $(U_j)$ together with the corresponding measures $\mu_j$, converges in $L^1(B)$, and let $u$ be the nonpositive plurisubharmonic representative of its limit. Then $u\in L^2(\mu)$ and
\begin{equation}\label{eq:canonical-grid-self-integral-convergence}
 \int_B|U_j-u|\,\mathrm d\mu\longrightarrow0,
 \qquad
 \int_B|U_j-u|\,\mathrm d\mu_j\longrightarrow0.
\end{equation}
In particular,
\[
 \int_B U_j(\ddc U_j)^n\longrightarrow\int_Bu\,\mathrm d\mu.
\]
\end{theorem}
\begin{proof}
If $m=0$, then $\mu=\mu_j=0$ and there is nothing to prove. We assume that $m>0$, and use the cubes and the polydiscs $P(z,r)$ of the proof of \cref{prop:canonical-grid-energy-control}; recall that $K$ denotes the support of $\mu$.

\emph{Step 1. Convergence against the fixed measure $\mu$.}
By \cref{thm:bounded-subsolution-local-energy-moments} with $p=2$ and the uniform bound for $e_2(U_j)$ in \cref{prop:canonical-grid-energy-control},
\[
 S=\sup_j\int_BU_j^2\,\mathrm d\mu<\infty.
\]
Fix $T>0$. The truncations $\max\{U_j,-T\}$ are plurisubharmonic with values in $[-T,0]$, and they converge to $\max\{u,-T\}$ in $L^1_{\loc}(B)$, since $s\mapsto\max\{s,-T\}$ is $1$-Lipschitz. The measure $\mu$ satisfies the hypothesis of \cref{lem:bounded-psh-fixed-measure-continuity}, because $\mu\le(\ddc w)^n$, so the truncations converge in $L^1(\mu)$. Since $|a^2-b^2|\le2T|a-b|$ for $a,b\in[-T,0]$, it follows that $\int\max\{u,-T\}^2\,\mathrm d\mu\le S$. Letting $T\to\infty$, monotone convergence gives $\int u^2\,\mathrm d\mu\le S$; in particular $u\in L^2(\mu)$, and $u$ is finite $\mu$-almost everywhere. For every measurable $v\le0$ we have the elementary tail bound
\[
 \int_B|v-\max\{v,-T\}|\,\mathrm d\mu
             \leq T^{-1}\int_Bv^2\,\mathrm d\mu,
\]
because $0\le\max\{v,-T\}-v\le v^2/T$. Applying it to $v=U_j$ and $v=u$, we get
\[
 \limsup_{j\to\infty}\int_B|U_j-u|\,\mathrm d\mu\le\frac{2S}T,
\]
and letting $T\to\infty$ proves the first convergence in \eqref{eq:canonical-grid-self-integral-convergence}.

\emph{Step 2. Cell averages.}
For $z\in K$ let $Q_j(z)$ be the cube of $\mathcal Q_j$ containing $z$, and write
\[
 \mathcal A_jv(z)=\frac1{|Q_j(z)|}\int_{Q_j(z)}v\,\mathrm dV
\]
for the average over this cube of a function $v$ integrable on it. If $v$ is $\mu_j$-integrable, then $v$ is integrable on every cube $Q$ with $\mu(Q)>0$; on a cube with $\mu(Q)=0$ on which $v$ is not integrable we set $\mathcal A_jv=0$, which does not affect integrals against $\mu$. Since $\mu$ is carried by $K$, the definition of $\mu_j$ gives
\begin{equation}\label{eq:canonical-grid-cell-pairing}
 \int_Bv\,\mathrm d\mu_j=\int_K\mathcal A_jv\,\mathrm d\mu.
\end{equation}
For a plurisubharmonic function $v\le0$ on $B$ and $z\in K$, the Cauchy--Schwarz inequality and the polydisc estimate in the proof of \cref{prop:canonical-grid-energy-control}, with $p=2$, give
\begin{equation}\label{eq:canonical-grid-cell-square}
 |\mathcal A_jv(z)|^2
 \leq\mathcal A_j(v^2)(z)\leq D_{n,2}|v(z)|^2;
\end{equation}
the right side is infinite if $v(z)=-\infty$. The polydisc estimate requires only $v(z)>-\infty$, so this applies also to the possibly unbounded function $u$.

Fix $z\in K$ with $u(z)>-\infty$ and $\varepsilon>0$, and put $c=u(z)+2\varepsilon$. By upper semicontinuity of $u$ and Hartogs' lemma (\cref{thm:psh-compactness-hartogs}), there is a neighborhood of $z$ on which $u\le c$ and $U_j\le c$ for all large $j$. Since the polydiscs $P(z,4\sqrt2\,h_j)$ shrink to $z$, for $v=u$ or $v=U_j$ and all large $j$ the function $v-c$ is plurisubharmonic and nonpositive on $P(z,4\sqrt2\,h_j)$. The polydisc estimate with $p=1$ then gives $\mathcal A_j(c-v)(z)\le D_{n,1}(c-v(z))$, while $\mathcal A_jv(z)\le c$ because $v\le c$ on $Q_j(z)$. Thus, for all large $j$,
\[
 c-D_{n,1}(c-v(z))\leq\mathcal A_jv(z)\leq c,
                   \qquad v=u\ \text{or}\ v=U_j.
\]
For $v=u$ this gives $u(z)+2\varepsilon(1-D_{n,1})\le\mathcal A_ju(z)\le u(z)+2\varepsilon$, so letting $j\to\infty$ and then $\varepsilon\downarrow0$ proves $\mathcal A_ju(z)\to u(z)$. The same argument shows that $\mathcal A_jU_j(z)\to u(z)$ whenever $U_j(z)\to u(z)$.

By the first convergence in \eqref{eq:canonical-grid-self-integral-convergence}, every subsequence of $(U_j)$ has a further subsequence converging to $u$ $\mu$-almost everywhere. Since $u$ is finite $\mu$-almost everywhere, it follows that $\mathcal A_jU_j\to u$ in $\mu$-measure, and $\mathcal A_ju\to u$ $\mu$-almost everywhere. By \eqref{eq:canonical-grid-cell-square} and Step~1, both sequences of cell averages are bounded in $L^2(\mu)$. A sequence $F_j$ which converges to zero in $\mu$-measure and is bounded in $L^2(\mu)$ converges to zero in $L^1(\mu)$: for $\varepsilon>0$ the Cauchy--Schwarz inequality gives
\[
 \int_K|F_j|\,\mathrm d\mu
 \leq\varepsilon m+\|F_j\|_{L^2(\mu)}
                 \mu(\{|F_j|>\varepsilon\})^{1/2}.
\]

Letting first $j\to\infty$ and then $\varepsilon\downarrow0$, and applying this to $F_j=\mathcal A_jU_j-u$ and $F_j=\mathcal A_ju-u$, we see that $\mathcal A_jU_j\to u$ and $\mathcal A_ju\to u$ in $L^1(\mu)$. The function $U_j-u$ is $\mu_j$-integrable, since $u$ is locally integrable and $\mu_j$ has a bounded density with compact support. Therefore \eqref{eq:canonical-grid-cell-pairing} gives
\begin{equation}\label{eq:canonical-grid-signed-integral}
 \int_B(U_j-u)\,\mathrm d\mu_j=\int_K(\mathcal A_jU_j-\mathcal A_ju)\,\mathrm d\mu\longrightarrow0.
\end{equation}

\emph{Step 3. The absolute difference.}
Put $g_j=\mathcal A_j((U_j-u)_+)$ on $K$. Let $z\in K$ with $u(z)>-\infty$, and let $c$ be as in Step~2. For all large $j$ we have $U_j\le c$ and $u\le c$ on $Q_j(z)$, so $(U_j-u)_+\le c-u$ there, and hence
\[
 0\leq g_j(z)\leq c-\mathcal A_ju(z).
\]
Since $\mathcal A_ju(z)\to u(z)$, letting $j\to\infty$ and then $\varepsilon\downarrow0$ shows that $g_j\to0$ $\mu$-almost everywhere. The Cauchy--Schwarz inequality and \eqref{eq:canonical-grid-cell-square} give
\[
 g_j(z)^2\leq2\mathcal A_j(U_j^2)(z)
                  +2\mathcal A_j(u^2)(z)
 \leq2D_{n,2}\bigl(U_j(z)^2+u(z)^2\bigr),
\]
so the $g_j$ are bounded in $L^2(\mu)$ by Step~1. The argument at the end of Step~2 gives $\int_Kg_j\,\mathrm d\mu\to0$, and by \eqref{eq:canonical-grid-cell-pairing} this integral equals $\int_B(U_j-u)_+\,\mathrm d\mu_j$. Combining this with \eqref{eq:canonical-grid-signed-integral} and the identity $|U_j-u|=2(U_j-u)_+-(U_j-u)$, we obtain the second convergence in \eqref{eq:canonical-grid-self-integral-convergence}.

Finally, since $(\ddc U_j)^n=\mu_j$,
\[
 \int_BU_j(\ddc U_j)^n=\int_B(U_j-u)\,\mathrm d\mu_j+\int_Bu\,\mathrm d\mu_j.
\]
The first term tends to zero by the second convergence in \eqref{eq:canonical-grid-self-integral-convergence}, and by \eqref{eq:canonical-grid-cell-pairing} and Step~2 the second term equals $\int_K\mathcal A_ju\,\mathrm d\mu$, which tends to $\int_Bu\,\mathrm d\mu$. This proves the last assertion.
\end{proof}
\index{local energy!canonical self-integral convergence}


\subsection{Capacity convergence and truncated measures}
\label{sec:canonical-grid-capacity-convergence}
The absolute energy error now replaces a common density bound in the comparison argument. Truncation controls the part of the measure carried by large negative values. All products in this argument still have bounded potentials.

For finite positive measures $\alpha,\beta$ on $B$ we use the total variation distance
\[
 \|\alpha-\beta\|_{\mathrm{TV}}
 =\sup\left\{\left|\int f\,\mathrm d(\alpha-\beta)\right|:
                   f\text{ Borel and real},\ |f|\leq1\right\}.
\]
By \cref{prop:continuous-density}, the same supremum is obtained with $f\in C_c(B)$: approximate in $L^1(\alpha+\beta)$ and clamp to $[-1,1]$.

\begin{lemma}\label{lem:zero-boundary-truncation-mass}
For $\varphi\in\mathcal E_0(B)$ and $T>0$ put $\varphi_T=\max\{\varphi,-T\}$. Then $\varphi_T\in\mathcal E_0(B)$ and, for every $p\geq1$,
\[
 \int_B(\ddc\varphi_T)^n=\int_B(\ddc\varphi)^n,
 \qquad e_p(\varphi_T)\leq e_p(\varphi),
\]
\begin{equation}\label{eq:zero-boundary-truncation-variation}
 \|(\ddc\varphi_T)^n-(\ddc\varphi)^n\|_{\mathrm{TV}}
                         \leq2e_p(\varphi)T^{-p}.
\end{equation}
For every Borel $E\subset B$,
\begin{equation}\label{eq:energy-measure-capacity-modulus}
 \int_E(\ddc\varphi)^n
       \leq T^n\operatorname{Cap}_B(E)+e_p(\varphi)T^{-p}.
\end{equation}
\end{lemma}
\begin{proof}
The two potentials agree on a boundary collar, so their difference has compact support in $B$. The truncated measure is finite: it agrees with the original measure on that collar, and is locally finite on the remaining compact set. Symmetry and the recursive definition give
\[
 (\ddc\varphi_T)^n-(\ddc\varphi)^n
 =\ddc\left((\varphi_T-\varphi)
       \sum_{k=0}^{n-1}(\ddc\varphi_T)^k
                              \wedge(\ddc\varphi)^{n-1-k}\right).
\]
The current in parentheses has compact support. Test with a smooth cutoff equal to one near that support; its $\ddc$ vanishes there. Integration proves equality of the total masses. The zero boundary limit is preserved, giving membership in $\mathcal E_0(B)$.

Maximum locality identifies the two measures on $\{\varphi>-T\}$. Equality of total mass therefore identifies their masses on $F_T=\{\varphi\leq-T\}$ as well. The truncated energy has weight $T^p$ on $F_T$, whereas the original weight is at least $T^p$. This proves the energy inequality. Both remaining measure masses are at most $e_p(\varphi)T^{-p}$, giving \eqref{eq:zero-boundary-truncation-variation}. Finally $(\varphi_T+T)/T$ is a capacity competitor. Split $E$ between $F_T$ and its complement, use maximum locality on the complement, and bound the mass on $F_T$ by its energy. This proves \eqref{eq:energy-measure-capacity-modulus}.
\end{proof}
\index{local energy!truncation mass and capacity}

\begin{proposition}\label{prop:energy-controlled-local-capacity-tails}
Let $\varphi_j\in\mathcal E_0(B)$ converge in $L^1_{\loc}(B)$ to a canonical nonpositive psh function $u$. Suppose $E_p=\sup_je_p(\varphi_j)<\infty$ for some $p\geq1$. For $T>0$,
\[
 \operatorname{Cap}_B(\{\varphi_j<-T\})
             \leq2^{n+p}(n+1)!E_pT^{-(n+p)}.
\]
For every $K\Subset B$ there is $C_{K,p}$, independent of $j,u,T$, such that
\begin{equation}\label{eq:energy-limit-capacity-tail}
 \operatorname{Cap}_B^*(K\cap\{u<-T\})
                \leq C_{K,p}E_p^{p/(n+p)}T^{-p}.
\end{equation}
In particular the pole set of $u$ has zero outer capacity locally.
\end{proposition}
\begin{proof}
Fix a capacity competitor $0\leq h\leq1$. For the first estimate restrict $(\ddc h)^n$ to $\{\varphi_j<-T\}$. This measure is finite because the sublevel set is relatively compact, and it is dominated by the measure of the bounded subsolution $h$. The tail estimate in \cref{thm:bounded-subsolution-local-energy-moments}, whose energy coefficient is independent of the restricted mass, gives the asserted bound. Take the supremum over $h$. Since the sublevel set is open, this also bounds its outer capacity.

For the limit estimate choose $K\subset U\Subset L^\circ\Subset B$, with $U$ open and $L$ compact. The bounded CLN estimate gives a uniform bound on $\int_L(\ddc h)^n$. Apply the moment estimate to the finite measure $\mathbf1_L(\ddc h)^n$ and the test $\varphi_j$. It gives
\[
 \int_L(-\varphi_j)^p(\ddc h)^n
                         \leq C_{K,p}E_p^{p/(n+p)}.
\]
For fixed $S>0$ the bounded truncations $\max\{\varphi_j,-S\}$ converge in local volume $L^1$ to $\max\{u,-S\}$. By \cref{lem:bounded-psh-fixed-measure-continuity} they converge in $L^1$ of the fixed restricted measure. The $p$th power is Lipschitz on $[0,S]$, so its integral also converges. Let $S\to\infty$ by monotone convergence. This gives the same bound for $\int_L(-u)^p(\ddc h)^n$. Chebyshev's inequality and the supremum over $h$ bound the capacity of the open set $U\cap\{u<-T\}$, which covers the set in \eqref{eq:energy-limit-capacity-tail}. Finally let $T\to\infty$ on each compact set to treat the poles.
\end{proof}

\begin{theorem}\label{thm:canonical-grid-capacity-convergence}
Under the hypotheses of \cref{thm:canonical-grid-self-integral-convergence}, the relabeled subsequence satisfies $U_j\to u$ locally in outer capacity. Put $E_2=\sup_je_2(U_j)$. Every bounded truncation $u_T=\max\{u,-T\}$ has finite total Monge--Amp\`ere mass at most $m$, and
\begin{equation}\label{eq:canonical-grid-truncated-measure-limit}
 \|(\ddc u_T)^n-\mu\|_{\mathrm{TV}}\leq2E_2T^{-2}.
\end{equation}
\end{theorem}
\begin{proof}
Write $a_j=\max\{U_j,-T\}$ and $b=\max\{u,-T\}$, with $T$ fixed for the moment. They lie in $[-T,0]$ and $a_j\to b$ in local volume $L^1$. Put
\[
 \varepsilon_j=\int_B|U_j-u|\,\mathrm d\mu_j\longrightarrow0.
\]
For $\delta>0$, the set $\{a_j<b-\delta/2\}$ is relatively compact: $b\leq0$ and $a_j$ has zero boundary limit. Apply \cref{lem:local-comparison-capacity-sublevels} with $s=t=\delta/2$, then use the truncation lemma:
\[
 \begin{aligned}
 \operatorname{Cap}_B(\{a_j<b-\delta\})
 &\leq2^{n+1}\delta^{-(n+1)}
                      \int_B(b-a_j)_+(\ddc a_j)^n\\
 &\leq2^{n+1}\delta^{-(n+1)}
                         \left(\varepsilon_j+\frac{2E_2}{T}\right).
 \end{aligned}
\]
Indeed truncation does not increase $|U_j-u|$, the positive integrand is at most $T$, and the total variation error is at most $2E_2/T^2$.

The upper difference has capacity tending to zero on every compact $L\Subset B$. Choose an open set $G$ of arbitrarily small capacity on whose complement $b$ is continuous. On $L\setminus G$, upper semicontinuity, Hartogs' lemma with a constant majorant and a finite cover give $a_j-b<\delta$ eventually. Thus $L\cap\{a_j>b+\delta\}\subset G$. Both bounded functions are quasi-continuous. With $K\Subset L^\circ$, \cref{lem:quasi-continuous-level-capacity} now turns the two ordinary level bounds into
\[
 \limsup_j\operatorname{Cap}_B^*(K\cap\{|a_j-b|>\delta\})
                      \leq C_n\delta^{-(n+1)}E_2/T.
\]
The preceding tail proposition controls $\{U_j<-T\}$ uniformly and $K\cap\{u<-T\}$. Outside those two sets the truncated and original differences agree. First let $j\to\infty$, then $T\to\infty$; all three capacity bounds tend to zero. This proves the stated local outer-capacity convergence, including the null pole set of $u$.

For each fixed $T$, truncation preserves this capacity convergence. Bounded-product continuity in \cref{thm:bt-capacity-continuity} gives
\[
 (\ddc a_j)^n\longrightarrow(\ddc u_T)^n
                                    \quad\text{weakly on }B.
\]
\Needspace{6\baselineskip}
The measures on the left have total mass $m$ by the truncation lemma. Testing with nonnegative compact cutoffs and exhausting $B$ proves that the limit has mass at most $m$. Also
\[
 \|(\ddc a_j)^n-\mu_j\|_{\mathrm{TV}}\leq2E_2/T^2.
\]
Pass to the limit against any real $f\in C_c(B)$ with $|f|\leq1$, using the weak convergence of $\mu_j$ to $\mu$. The resulting bound extends to bounded Borel tests by \cref{prop:continuous-density} for the finite measure $(\ddc u_T)^n+\mu$. Taking their supremum proves \eqref{eq:canonical-grid-truncated-measure-limit}.
\end{proof}
\index{convergence!capacity for energy-controlled potentials}
\index{Monge--Amp\`ere measure!truncated canonical limit}


\subsection{A bounded subsolution and the limiting equation}
\label{sec:bounded-subsolution-ball-equation}
We can now identify the canonical limit. A bounded local subsolution supplies a zero-boundary barrier near the support of the prescribed measure. Comparison with a strict perturbation of that barrier gives a bound for the limit and recovers its boundary values.

\begin{lemma}\label{lem:capacity-controlled-measure-weights}
Let $\nu_j,\nu$ be finite positive measures supported in a fixed compact set $K_0\Subset B$, with uniformly bounded masses, and suppose $\nu_j\to\nu$ weakly. Assume there is a nondecreasing function $\eta:[0,+\infty]\to[0,+\infty]$ with $\eta(t)\to0$ as $t\downarrow0$ such that, for every Borel $E\subset B$,
\[
 \nu_j(E),\ \nu(E)\leq\eta(\operatorname{Cap}_B^*(E)).
\]
Let $f_j,f$ be real Borel functions with a common finite bound, let $f$ be quasi-continuous near $K_0$, and suppose $f_j\to f$ in outer capacity on $K_0$. Then
\[
 \int_Bf_j\,\mathrm d\nu_j\longrightarrow\int_Bf\,\mathrm d\nu.
\]
For every $c\in\mathbb R$ with $\nu(\{f=c\})=0$, also
\begin{equation}\label{eq:capacity-controlled-moving-level-sets}
 \nu_j(\{f_j<c\})\longrightarrow\nu(\{f<c\}).
\end{equation}
\end{lemma}
\begin{proof}
Write $C$ for the common function bound and $M$ for the mass bound. First keep the weight $f$ fixed. On a compact neighborhood of $K_0$, remove an open $G$ of arbitrarily small capacity so that $f$ is continuous on the remaining compact set. Its continuous extension from \cref{sec:bt-quasicontinuity}, multiplied by a compact cutoff equal to one near $K_0$, gives $g\in C_c(B)$ with $|g|\leq C$ and $g=f$ on $K_0\setminus G$. Weak convergence applies to $g$, and the two replacement errors are at most $2C\nu_j(G)$ and $2C\nu(G)$. Let the capacity of $G$ tend to zero. This proves convergence for the fixed quasi-continuous weight.

For the varying weights and $\delta>0$,
\[
 \int_B|f_j-f|\,\mathrm d\nu_j
 \leq\delta M+2C\eta\left(
       \operatorname{Cap}_B^*(K_0\cap\{|f_j-f|>\delta\})\right).
\]
Let $j\to\infty$, then $\delta\downarrow0$, and combine with the fixed-weight result.

For the level-set assertion choose continuous piecewise linear functions $\sigma^-_s,\sigma^+_s$ with values in $[0,1]$ such that
\[
 \mathbf1_{(-\infty,c-s)}\leq\sigma^-_s
 \leq\mathbf1_{(-\infty,c)}\leq\sigma^+_s
 \leq\mathbf1_{(-\infty,c+s]},\qquad s>0.
\]
For example the lower ramp falls linearly from one at $c-s$ to zero at $c$, and the upper one from one at $c$ to zero at $c+s$. Their compositions with $f_j$ converge in capacity to the corresponding compositions with $f$, whose quasi-continuity is preserved. Apply the first assertion to these compositions and squeeze. As $s\downarrow0$, dominated convergence and the null level set give \eqref{eq:capacity-controlled-moving-level-sets}.
\end{proof}

\begin{corollary}\label{cor:energy-controlled-fixed-measure-limit}
Suppose $\varphi_j\in\mathcal E_0(B)$, $\sup_je_2(\varphi_j)<\infty$, and $\varphi_j\to u$ in local volume $L^1$, with the canonical nonpositive psh limit. If $\nu$ is a finite positive Borel measure and $\nu\leq(\ddc w)^n$ for a bounded $w\in\PSH(B)$, then $u\in L^2(\nu)$ and $\varphi_j\to u$ in $L^1(\nu)$.
\end{corollary}
\begin{proof}
The moment theorem with $p=2$ gives uniform $L^2(\nu)$ bounds for $\varphi_j$. For fixed $T$, the bounded truncations converge in $L^1(\nu)$ by \cref{lem:bounded-psh-fixed-measure-continuity}, and their squares converge in integral. Letting $T\to\infty$ gives the same $L^2$ bound for $u$. The error between a negative function and its truncation has $L^1(\nu)$ norm at most $T^{-1}$ times its squared $L^2(\nu)$ norm. Apply this to $\varphi_j,u$, first taking $j\to\infty$ at fixed $T$ and then $T\to\infty$.
\end{proof}

\begin{theorem}\label{thm:compact-bounded-subsolution-ball-dirichlet}
Let $\mu$ be a finite positive measure of compact support in a ball $B\subset\mathbb C^n$, $n\geq1$. If $\mu\leq(\ddc w)^n$ for a bounded $w\in\PSH(B)$, there is a unique bounded psh function $U$ with
\[
 (\ddc U)^n=\mu,\qquad
 \lim_{B\ni z\to\zeta}U(z)=0\quad(\zeta\in\partial B).
\]
It belongs to $\mathcal E_0(B)$ and is the pointwise envelope of bounded psh subsolutions with upper boundary limits at most zero. The entire canonical sequence in \cref{prop:canonical-grid-energy-control} converges to $U$ in $L^1(B)$ and locally in outer capacity, and satisfies
\[
 \int_B|U_j-U|\,\mathrm d\mu
       +\int_B|U_j-U|\,\mathrm d\mu_j\longrightarrow0.
\]
\end{theorem}
\begin{proof}
If $\mu=0$, the homogeneous zero-boundary solution and all its canonical approximants are zero. The maximum principle identifies the stated envelope. Write $m=\mu(B)>0$ for the remaining case.

\emph{Step 1. A zero-boundary barrier.}
Write $B=B(a,R)$, $\rho(z)=|z-a|^2-R^2$ and $w_0=w-\sup_Bw-1$. Then $-M-1\leq w_0\leq-1$, where $M$ is the oscillation of $w$. Choose $A>0$ so large that $A\rho<-M-2$ on a neighborhood of $\operatorname{supp}\mu$. The function
\[
 W=\max\{w_0,A\rho\}
\]
equals $w_0$ on that neighborhood and equals $A\rho$ on a boundary collar. It is bounded, negative and has zero boundary values; maximum locality gives $(\ddc W)^n\geq\mu$.

Its total mass is finite, since the collar measure is smooth and the remaining interior measure is locally finite. Apply the compact-support polynomial-current identity in the proof of \cref{lem:zero-boundary-truncation-mass} to $W$ and $A\rho$. Their difference has compact support, so testing its $\ddc$ with a cutoff equal to one there gives
\[
 \int_B(\ddc W)^n=A^n\int_B(\ddc\rho)^n.
\]
The same argument gives finite mass for $W_\varepsilon=W+\varepsilon\rho$, $\varepsilon>0$, whose collar is $(A+\varepsilon)\rho$. Positivity of mixed products gives
\begin{equation}\label{eq:local-subsolution-strict-barrier}
 \nu_\varepsilon:=(\ddc W_\varepsilon)^n
                       \geq\mu+\varepsilon^n(\ddc\rho)^n.
\end{equation}

\emph{Step 2. Comparison along a canonical subsequence.}
Choose an $L^1(B)$-convergent subsequence of $U_j$, with limit $u$. The preceding results give local outer-capacity convergence and a uniform second energy bound $E_2$. By the corollary, $U_j\to u$ in $L^1(\nu_\varepsilon)$ for each fixed $\varepsilon>0$.

The grid measures have supports in a fixed $K_0\Subset B$. They have a common capacity modulus by \eqref{eq:energy-measure-capacity-modulus}:
\[
 \mu_j(E)\leq T^n\operatorname{Cap}_B^*(E)+E_2T^{-2}.
\]
The original measure satisfies $\mu(E)\leq M^n\operatorname{Cap}_B^*(E)$. Define $\eta(t)$ as the supremum of their masses on Borel sets of outer capacity at most $t$. First choosing $T$ large and then making the capacity small shows that this common nondecreasing modulus tends to zero. Thus the weight lemma applies to $\mu_j,\mu$.

Fix $\varepsilon>0$. Choose $0<\delta<1$ such that the level set $\{u-W_\varepsilon=-\delta\}$ has zero mass for both $\mu$ and $\nu_\varepsilon$. Only countably many levels are excluded, since both measures are finite. For $T>\|W_\varepsilon\|_\infty+2$, the inequalities
\[
 U_j<W_\varepsilon-\delta,\qquad u<W_\varepsilon-\delta
\]
are unchanged when $U_j,u$ are replaced by their truncations at $-T$. The bounded differences between those truncations and $W_\varepsilon$ are quasi-continuous and converge in capacity. Hence \eqref{eq:capacity-controlled-moving-level-sets} gives
\[
 \mu_j(\{U_j<W_\varepsilon-\delta\})
                       \longrightarrow\mu(\{u<W_\varepsilon-\delta\}).
\]
The fixed-measure $L^1(\nu_\varepsilon)$ convergence gives the corresponding convergence for $\nu_\varepsilon$. To see this directly, outside $\{|U_j-u|>s\}$ the two sublevel indicators can differ only in $\{|u-W_\varepsilon+\delta|\leq s\}$. The first set has mass tending to zero; let $s\downarrow0$ in the second, using the null level set.

The comparison set for $U_j$ and $W_\varepsilon-\delta$ is relatively compact because $U_j$ has zero boundary limit and $W_\varepsilon-\delta\leq-\delta$. Bounded comparison gives
\[
 \nu_\varepsilon(\{U_j<W_\varepsilon-\delta\})
                       \leq\mu_j(\{U_j<W_\varepsilon-\delta\}).
\]
Pass to the limit. By \eqref{eq:local-subsolution-strict-barrier}, the set $\{u<W_\varepsilon-\delta\}$ has zero Euclidean volume. Thus $u\geq W_\varepsilon-\delta$ almost everywhere, and canonical psh representatives preserve this ordering at every point. Let $\delta\downarrow0$ through admissible levels, then $\varepsilon\downarrow0$. We obtain $W\leq u\leq0$.

The barrier makes $u$ bounded and gives its zero boundary limits. For every $T>\|u\|_\infty$, its truncation is $u$ itself. The estimate \eqref{eq:canonical-grid-truncated-measure-limit}, followed by $T\to\infty$, proves $(\ddc u)^n=\mu$. Its mass is $m$, so $u\in\mathcal E_0(B)$.

\emph{Step 3. Uniqueness and the full sequence.}
Let $v$ be a bounded psh subsolution with upper boundary limits at most zero, and put $v_\varepsilon=v+\varepsilon\rho$. Then $(\ddc v_\varepsilon)^n\geq\mu+\varepsilon^n(\ddc\rho)^n$. Comparison with $u$ gives
\[
 \mu(\{u<v_\varepsilon\})
       +\varepsilon^n\int_{\{u<v_\varepsilon\}}(\ddc\rho)^n
          \leq\mu(\{u<v_\varepsilon\}).
\]
The two $\mu$ integrals are finite. Their cancellation makes the strict set volume-null, so canonical representatives give $u\geq v_\varepsilon$ everywhere. Let $\varepsilon\downarrow0$. Since $u$ is itself a subsolution, this proves the pointwise envelope assertion. Interchanging two bounded solutions proves uniqueness; call the solution $U$.

Every subsequence of the canonical sequence has a further $L^1(B)$-convergent subsequence. The argument above identifies every such limit with $U$; \cref{thm:canonical-grid-self-integral-convergence,thm:canonical-grid-capacity-convergence} give its energy and capacity convergences. If any asserted full-sequence convergence failed, a subsequence separated from its stated limit would have such a further convergent subsequence, a contradiction. This proves all the assertions.
\end{proof}
\index{Dirichlet problem!bounded subsolution and compact measure}


\subsection{Mixed energy and truncation}
\label{sec:mixed-local-energy}
Canonical approximants have controlled energies but need not have a common bound as functions. Mixed products must therefore be compared with their bounded truncations before continuity can be used. A mixed energy estimate makes the discarded sets uniformly small for both products.

\begin{theorem}\label{thm:bounded-mixed-local-energy}
For $u_0,\ldots,u_n\in\mathcal E_0(B)$,
\begin{equation}\label{eq:bounded-mixed-local-energy}
 \int_B(-u_0)\,\ddc u_1\wedge\cdots\wedge\ddc u_n
 \leq C_n\prod_{l=0}^n e_1(u_l)^{1/(n+1)},
 \qquad C_n=2^{n+2}.
\end{equation}
In particular the integral is finite, also when some energy is zero.
\end{theorem}
\begin{proof}
Write $B=B(a,R)$ and $\rho=|z-a|^2-R^2$. We first prove the inequality with constant one for potentials which equal positive multiples of $\rho$ on boundary collars. We then approximate arbitrary members of $\mathcal E_0(B)$ by such potentials, retaining an explicit bound for their energies.

\emph{Step 1. The collar case.}
For this step assume that each $u_l=c_l\rho$ near $\partial B$, with $c_l>0$. All mixed products have finite total mass: their measures are smooth on a common collar and locally finite in the remaining compact interior. Put
\[
 I(u_0,\ldots,u_n)
       =\int_B(-u_0)\,\ddc u_1\wedge\cdots\wedge\ddc u_n.
\]
For smooth potentials with zero boundary values and
$S=\ddc u_2\wedge\cdots\wedge\ddc u_n$, Stokes' formula gives
\[
 I(u_0,u_1,u_2,\ldots,u_n)
          =\int_B\mathrm du_0\wedge\mathrm d^cu_1\wedge S.
\]
The boundary term vanishes because $u_0=0$ there. In dimension one $S=1$. Since $S$ is positive of type $(n-1,n-1)$, the real bilinear form
\[
 (b,c)\longmapsto\int_B\mathrm db\wedge\mathrm d^cc\wedge S
\]
is symmetric and nonnegative on the diagonal. Indeed its $(1,1)$ factor is
$(\mathrm i/(4\pi))(\partial b\wedge\bar\partial c+
\partial c\wedge\bar\partial b)$, and positivity applies when $b=c$. Applying nonnegativity to $b+tc$, $t\in\mathbb R$, proves Cauchy--Schwarz, including a zero diagonal term. Consequently
\begin{equation}\label{eq:mixed-local-energy-two-slot}
 I(b,c,u_2,\ldots,u_n)^2
 \leq I(b,b,u_2,\ldots,u_n)I(c,c,u_2,\ldots,u_n).
\end{equation}
The same symmetry interchanges the first two slots. Symmetry of the Hessian factors interchanges the remaining slots, so $I$ is symmetric in all $n+1$ slots.

These assertions also hold for bounded collar potentials. Extend each $u_l$ across the boundary using $c_l\rho$ and convolve with the same smooth radial kernel of radius $\varepsilon$. For a constant $\kappa>0$ determined by that kernel,
\[
 u_{l,\varepsilon}=u_l*\rho_\varepsilon-c_l\kappa\varepsilon^2
\]
equals $c_l\rho$ on a smaller fixed collar. It is smooth and psh on a neighborhood of $\overline B$, has zero boundary values, and is nonpositive by the maximum principle. The raw radial convolutions decrease to $u_l$. Their weighted products converge on compact tests by \cref{thm:bt-decreasing-continuity}; the subtracted constants change weighted integrals by terms tending to zero, by the local mass bounds. On the fixed collar the corrected potentials and their products already equal their limiting quadratic expressions. A compact cutoff for the interior therefore gives convergence of all the displayed global integrals. Passing to the limit proves full symmetry and \eqref{eq:mixed-local-energy-two-slot}. The corrected smoothings themselves need not decrease.

Every $e_1(u_l)$ is positive, because its collar integral is positive. Multiply $u_l$ by $e_1(u_l)^{-1/(n+1)}$ to normalize each energy to one. Consider the finitely many ordered $(n+1)$-tuples whose entries are drawn from these normalized potentials. Let $F$ be the largest value of $I$ on those tuples, and among the maximizing tuples choose one for which the sum of the squares of the entry multiplicities is largest. Diagonal tuples have value one, so $F\geq1$.

If the chosen tuple contains two different entries $b,c$, symmetry puts them in the first two slots. Both quantities on the right of \eqref{eq:mixed-local-energy-two-slot} are at most $F$. Since its left side is $F^2$, both replacements also maximize $I$. If $b,c$ have multiplicities $r,s$, replacing one $c$ by $b$ changes the sum of squares by $2(r-s)+2$; the opposite replacement changes it by $2(s-r)+2$. At least one change is positive, contradicting the choice of tuple. Thus the maximizing tuple is diagonal and $F=1$. Undoing the normalization proves the constant-one estimate in the collar case.

\emph{Step 2. Approximation with controlled energy.}
Fix $u\in\mathcal E_0(B)$ and let
\[
 \delta_j=j^{-n-1},\qquad v_j=\max\{u-\delta_j,j\rho\},
 \qquad m_\rho=\int_B(\ddc\rho)^n.
\]
The positive shift below $u$ makes $v_j=j\rho$ on a boundary collar: $u-\delta_j\leq-\delta_j$, whereas $j\rho> -\delta_j$ sufficiently near the boundary. Thus $v_j\in\mathcal E_0(B)$ and the compact-support current identity gives
\[
 \int_B(\ddc v_j)^n=j^n m_\rho.
\]
Moreover $v_j\to u$ uniformly on $B$. To check this, fix $s>0$. On a sufficiently thin boundary collar, $-s<u\leq0$ and $u-\delta_j\leq v_j\leq0$. On its compact complement, $j\rho<u-\delta_j$ for all large $j$, because $u$ is bounded. The uniform error is therefore at most $s+\delta_j$.

For $t>4\delta_j$ the two inclusions
\[
 \{v_j<-t\}\subset\{u<v_j/2-t/4\}
                       \subset\{u<-t/4\}
\]
follow from $u\leq v_j+\delta_j$ and $v_j\leq0$. The middle set is relatively compact, since the boundary limits of $u$ and $v_j/2-t/4$ are $0$ and $-t/4$. Bounded comparison gives
\[
 \int_{\{v_j<-t\}}(\ddc v_j)^n
       \leq2^n\int_{\{u<-t/4\}}(\ddc u)^n.
\]
Layer integration, using the total mass bound for $0<t\leq4\delta_j$, yields
\begin{equation}\label{eq:quadratic-collar-energy-approximation}
 e_1(v_j)\leq4\delta_j j^n m_\rho+2^{n+2}e_1(u)
                 =4m_\rho/j+C_ne_1(u).
\end{equation}

Apply this construction to every $u_l$. Their approximants are uniformly bounded, converge uniformly, and hence converge locally in capacity to their possibly discontinuous limits. Weighted capacity continuity, \cref{thm:bt-capacity-continuity}, gives convergence of $I$ on compact cutoffs with values in $[0,1]$. The collar estimate and \eqref{eq:quadratic-collar-energy-approximation} bound those tested integrals by
\[
 \prod_{l=0}^n(4m_\rho/j+C_ne_1(u_l))^{1/(n+1)}.
\]
Let $j\to\infty$ and exhaust $B$ by nonnegative compact cutoffs increasing to one. This proves \eqref{eq:bounded-mixed-local-energy}. The same limiting argument covers zero energies without normalizing them.
\end{proof}
\index{local energy!mixed inequality}

\Needspace{12\baselineskip}
\begin{corollary}\label{cor:mixed-products-energy-capacity-continuity}
For $1\leq l\leq n$ let $u_{l,j}\in\mathcal E_0(B)$, with $e_1(u_{l,j})\leq E<\infty$. Suppose $u_{l,j}\to u_l$ locally in capacity, where each $u_l$ is bounded, nonpositive and psh. Then
\[
 \ddc u_{1,j}\wedge\cdots\wedge\ddc u_{n,j}
          \longrightarrow\ddc u_1\wedge\cdots\wedge\ddc u_n
\]
weakly on $B$.
\end{corollary}
\begin{proof}
For $T>0$ write $a_{l,j,T}=\max\{u_{l,j},-T\}$, and denote the original and truncated mixed measures by $R_j$ and $R_{j,T}$. The truncation lemma gives $e_1(a_{l,j,T})\leq E$. Apply the theorem with extra weight $-u_{l,j}$ to $R_j$, and with extra weight $-a_{l,j,T}$ to $R_{j,T}$. Each resulting integral is at most $C_nE$.

Let $D_{j,T}=\bigcup_{l=1}^n\{u_{l,j}\leq-T\}$. On its $l$-th constituent both extra weights are at least $T$, so
\[
 R_j(D_{j,T}),\ R_{j,T}(D_{j,T})\leq nC_nE/T.
\]
Iterated strict maximum locality gives $R_j=R_{j,T}$ on the complement, where every $u_{l,j}>-T$. Thus for each compactly supported continuous $\chi$,
\[
 \left|\int_B\chi\,(R_j-R_{j,T})\right|
                    \leq2nC_nE\|\chi\|_\infty/T.
\]
The tie sets were included in $D_{j,T}$. For fixed $T>\max_l\|u_l\|_\infty$, the truncated potentials converge locally in capacity to $u_l$ and have a common bound. Bounded capacity continuity identifies the weak limit of $R_{j,T}$ with the stated limiting mixed product. First let $j\to\infty$ at fixed $T$, then let $T\to\infty$ in the error estimate.
\end{proof}
\index{Bedford--Taylor product!energy-controlled continuity}


\section{Mixed products for a common measure}
\label{sec:common-measure-mixed-products}
\Cref{cor:bounded-multifactor-continuous-density} bounds a mixed Monge--Amp\`ere product from below when the pure Monge--Amp\`ere measures of the factors have continuous lower densities with respect to Lebesgue measure. For non-pluripolar products in \cref{chap:q09} we need the same inequality relative to an arbitrary positive Radon measure $\mu$, which may be singular with respect to Lebesgue measure. Instead of regularizing the potentials by convolution, as in the proof of that corollary, we restrict each pure Monge--Amp\`ere measure to a compact set and replace it by its cube averages, as in \cref{prop:canonical-grid-energy-control}. The averaged measures have piecewise constant densities, so the case of Lebesgue measure applies on each cube, and H\"older's inequality compares the resulting bound with the cube averages of the expected lower bound. The convergence results of \cref{sec:local-energy-moments} then allow us to pass to the limit.

\begin{theorem}\label{thm:bounded-multifactor-common-measure}
Let $\Omega\subset\mathbb C^n$ be open, $n\geq1$, and let $u_1,\ldots,u_n$ be locally bounded plurisubharmonic functions on $\Omega$. Let $\mu$ be a positive Radon measure on $\Omega$, and let $f_1,\ldots,f_n$ be nonnegative Borel functions on $\Omega$. If
\[
 (\ddc u_i)^n\geq f_i\mu\qquad(1\leq i\leq n),
\]
then
\begin{equation}\label{eq:bounded-multifactor-common-measure}
 \ddc u_1\wedge\cdots\wedge\ddc u_n
                  \geq\left(\prod_{i=1}^nf_i^{1/n}\right)\mu.
\end{equation}
For $n=1$ the assertion is the hypothesis.
\end{theorem}
\begin{proof}
Let $n\geq2$. We use the generalized H\"older inequality
\begin{equation}\label{eq:generalized-holder-geometric-mean}
 \int_A\prod_{i=1}^ng_i^{1/n}\,\mathrm d\nu\leq\prod_{i=1}^n\Bigl(\int_Ag_i\,\mathrm d\nu\Bigr)^{1/n}
\end{equation}
for a positive measure $\nu$, a measurable set $A$ and nonnegative measurable functions $g_i$. If one of the integrals $c_i=\int_Ag_i\,\mathrm d\nu$ vanishes, then $g_i=0$ $\nu$-almost everywhere on $A$ and both sides vanish, with the convention $0\cdot\infty=0$; if none vanishes and one is infinite, the right side is infinite. Otherwise, integrating the arithmetic--geometric mean inequality $\prod_i(g_i/c_i)^{1/n}\leq\frac1n\sum_ig_i/c_i$ over $A$ gives $\int_A\prod_i(g_i/c_i)^{1/n}\,\mathrm d\nu\leq1$, which is \eqref{eq:generalized-holder-geometric-mean}.\index{H\"older inequality!generalized} For every compact $K\subset\Omega$ we have $\int_Kf_i\,\mathrm d\mu\leq(\ddc u_i)^n(K)<\infty$ by hypothesis, so \eqref{eq:generalized-holder-geometric-mean} shows that the right side of \eqref{eq:bounded-multifactor-common-measure} is a positive Radon measure. We first prove the theorem when $\Omega$ is a ball $B=B(a_0,R)$ and all $u_i$ belong to $\mathcal E_0(B)$ (Steps 1 and 2), and then reduce the general case to this one by modifying the potentials near the boundary of a ball (Step 3). Put $\rho(z)=|z-a_0|^2-R^2$.

\emph{Step 1. Exhaustion by compactly supported measures.}
Choose closed balls $K_l\subset B$ concentric with $B$, increasing with $l$ and exhausting $B$, and set
\[
 \alpha_{i,l}=\mathbf1_{K_l}(\ddc u_i)^n.
\]
These measures are finite since $u_i\in\mathcal E_0(B)$, they have compact support in $B$, and they are dominated by the Monge--Amp\`ere measure of the bounded plurisubharmonic function $u_i$. By \cref{thm:compact-bounded-subsolution-ball-dirichlet}, the equation $(\ddc A)^n=\alpha_{i,l}$ has a bounded solution $A_{i,l}\in\mathcal E_0(B)$ with zero boundary values, and $A_{i,l}$ is the largest bounded plurisubharmonic function $v$ on $B$ with $(\ddc v)^n\geq\alpha_{i,l}$ and upper boundary limits at most zero. The function $u_i$ has these properties for every $l$, and so has $A_{i,l+1}$, because $\alpha_{i,l+1}\geq\alpha_{i,l}$. Therefore
\[
 u_i\leq A_{i,l+1}\leq A_{i,l}\leq0.
\]

We claim that $A_{i,l}\downarrow u_i$ as $l\to\infty$. The decreasing limit $A_{i,\infty}$ satisfies $u_i\leq A_{i,\infty}\leq0$, so it is a bounded plurisubharmonic function (\cref{thm:psh-decreasing-nets}) with zero boundary values. By \cref{thm:bt-decreasing-continuity}, $(\ddc A_{i,\infty})^n$ is the weak limit of $\alpha_{i,l}$, that is,
\[
 (\ddc A_{i,\infty})^n=(\ddc u_i)^n,
\]
and this measure has finite total mass. Let $\varepsilon>0$ and $S=\{u_i<A_{i,\infty}+\varepsilon\rho\}$. Since $u_i-(A_{i,\infty}+\varepsilon\rho)$ tends to zero at $\partial B$, the comparison theorem (\cref{thm:bt-comparison}) applies. Expanding $(\ddc(A_{i,\infty}+\varepsilon\rho))^n$ and dropping the mixed terms, which are positive, we obtain
\[
 \int_S(\ddc u_i)^n+\varepsilon^n\int_S(\ddc\rho)^n
 \leq\int_S\bigl(\ddc(A_{i,\infty}+\varepsilon\rho)\bigr)^n
 \leq\int_S(\ddc u_i)^n.
\]
The integrals of $(\ddc u_i)^n$ are finite, so $S$ is Lebesgue null. Hence $u_i\geq A_{i,\infty}+\varepsilon\rho$ almost everywhere, and therefore everywhere on $B$, by the uniqueness assertion of \cref{thm:psh-distribution-representative} applied to $\max\{u_i,A_{i,\infty}+\varepsilon\rho\}$ and $u_i$. Letting $\varepsilon\downarrow0$ gives $A_{i,\infty}=u_i$, as claimed.

\emph{Step 2. A common grid.}
Fix $l$. Choose a sequence of partitions of $\mathbb C^n=\mathbb R^{2n}$ into half-open cubes whose side lengths decrease to zero, the same for all $i$. Let $\alpha_{i,l,j}$ be the corresponding cube averages of $\alpha_{i,l}$, as in \cref{prop:canonical-grid-energy-control}, and let $A_{i,l,j}\in\mathcal E_0(B)$ be the corresponding canonical approximants of $\alpha_{i,l}$, that is, the zero-boundary solutions of $(\ddc A_{i,l,j})^n=\alpha_{i,l,j}$; they are defined for all large $j$. Discarding finitely many indices $j$, we may assume that they are defined for all $j$ and all $i$, and that every cube meeting $K_l$ has its closure in $B$. By \cref{thm:compact-bounded-subsolution-ball-dirichlet}, $A_{i,l,j}\to A_{i,l}$ locally in capacity as $j\to\infty$. By \cref{prop:canonical-grid-energy-control}, applied with the bounded subsolution $u_i$, the energies $e_1(A_{i,l,j})$ are bounded uniformly in $j$; since there are only finitely many indices $i$, they have a common bound. (If $\alpha_{i,l}=0$, all these solutions vanish.) Therefore \cref{cor:mixed-products-energy-capacity-continuity} gives
\begin{equation}\label{eq:common-measure-grid-mixed-limit}
 \bigwedge_{i=1}^n\ddc A_{i,l,j}
                      \longrightarrow\bigwedge_{i=1}^n\ddc A_{i,l}
\end{equation}
weakly on $B$.

Let $\lambda_l=\mathbf1_{K_l}\prod_if_i^{1/n}\mu$, a finite measure, and let $\lambda_{l,j}$ be its cube averages. For every cube $Q$, the hypothesis gives $\alpha_{i,l}(Q)\geq\int_{Q\cap K_l}f_i\,\mathrm d\mu$, and \eqref{eq:generalized-holder-geometric-mean} then gives
\begin{equation}\label{eq:common-measure-grid-holder}
 \prod_{i=1}^n\left(\frac{\alpha_{i,l}(Q)}{|Q|}\right)^{1/n}
 \geq\frac1{|Q|}\int_{Q\cap K_l}\prod_{i=1}^nf_i^{1/n}\,\mathrm d\mu
 =\frac{\lambda_l(Q)}{|Q|}.
\end{equation}
Let $Q$ be a cube with $\lambda_l(Q)>0$; it meets $K_l$, so its interior lies in $B$. There the measures $(\ddc A_{i,l,j})^n=\alpha_{i,l,j}$ have the constant densities $\alpha_{i,l}(Q)/|Q|$ with respect to Lebesgue measure $\mathrm dV$. Recall that $\eta^n=c_n\,\mathrm dV$ with $\eta=\ddc|z|^2$ and $c_n=n!\pi^{-n}$, by \eqref{eq:bt-determinant-normalization}. Applying \cref{cor:bounded-multifactor-continuous-density} on the open set $\Int Q$ to the constant densities $c_n^{-1}\alpha_{i,l}(Q)/|Q|$ with respect to $\eta^n$, and using \eqref{eq:common-measure-grid-holder}, we obtain on $\Int Q$
\[
 \bigwedge_{i=1}^n\ddc A_{i,l,j}
 \geq\prod_{i=1}^n\left(\frac{\alpha_{i,l}(Q)}{c_n|Q|}\right)^{1/n}\eta^n
 =\prod_{i=1}^n\left(\frac{\alpha_{i,l}(Q)}{|Q|}\right)^{1/n}\mathrm dV
 \geq\lambda_{l,j}.
\]
The factor $c_n$ cancels because the $n$ exponents $1/n$ sum to one. The measure $\lambda_{l,j}$ vanishes on the cubes with $\lambda_l(Q)=0$ and does not charge the cube boundaries, which are Lebesgue null. Since the mixed product is positive, we conclude that
\[
 \bigwedge_{i=1}^n\ddc A_{i,l,j}\geq\lambda_{l,j}\quad\text{on }B.
\]

As $j\to\infty$, $\lambda_{l,j}\to\lambda_l$ weakly. Indeed, only finitely many cubes have positive $\lambda_l$-mass, and for $\chi\in C_c(B)$, replacing the values of $\chi$ by its cube averages changes $\int\chi\,\mathrm d\lambda_l$ by at most $\lambda_l(B)$ times the modulus of continuity of $\chi$ at the cube diameter. Letting $j\to\infty$ and using \eqref{eq:common-measure-grid-mixed-limit}, we obtain $\bigwedge_i\ddc A_{i,l}\geq\lambda_l$. Finally let $l\to\infty$. By Step 1, $A_{i,l}\downarrow u_i$ with $u_i\leq A_{i,l}\leq0$, so \cref{thm:bt-decreasing-continuity} gives $\bigwedge_i\ddc A_{i,l}\to\bigwedge_i\ddc u_i$ weakly, while $\lambda_l$ increases to $\prod_if_i^{1/n}\mu$. Testing with nonnegative functions in $C_c(B)$ and using monotone convergence on the right, we obtain \eqref{eq:bounded-multifactor-common-measure} for potentials in $\mathcal E_0(B)$.

\emph{Step 3. Localization.}
Now let $u_1,\ldots,u_n$ be locally bounded on $\Omega$. Choose a ball $B=B(a_0,R)$ with $\overline B\subset\Omega$ and a compact set $K\subset B$. Since adding constants does not change the products, we may assume that $u_i\leq-1$ on $B$. With $\rho(z)=|z-a_0|^2-R^2$ and constants $C_i>0$ so large that $C_i\rho<\inf_Bu_i$ on a neighborhood of $K$, put
\[
 \widetilde u_i=\max\{u_i,C_i\rho\}.
\]
Then $\widetilde u_i=u_i$ on a neighborhood of $K$, and $\widetilde u_i=C_i\rho$ near $\partial B$, because $C_i\rho\to0$ at $\partial B$ while $u_i\leq-1$. Thus $\widetilde u_i\in\mathcal E_0(B)$: its Monge--Amp\`ere measure is smooth near $\partial B$ and locally finite on the rest of $B$. By locality of the products on the open neighborhood of $K$ where $\widetilde u_i=u_i$, we have $(\ddc\widetilde u_i)^n\geq f_i\mathbf1_K\mu$ on $B$. The case already proved, applied to the measure $\mathbf1_K\mu$, gives $\bigwedge_i\ddc\widetilde u_i\geq\prod_if_i^{1/n}\mathbf1_K\mu$, and on a neighborhood of $K$ the left side equals $\bigwedge_i\ddc u_i$ by locality. Thus $P=\bigwedge_i\ddc u_i$ and $\lambda=\prod_if_i^{1/n}\mu$ satisfy $P(E)\geq\lambda(E)$ for every Borel set $E\subset K$. The interiors of such compact sets $K$ cover $\Omega$, and since $\Omega$ has a countable basis, countably many of them, say $\Int K_1,\Int K_2,\ldots$, already cover $\Omega$. A Borel set $E\subset\Omega$ is the disjoint union of the Borel sets $E_m=(E\cap K_m)\setminus\bigcup_{m'<m}K_{m'}$, and therefore $P(E)=\sum_mP(E_m)\geq\sum_m\lambda(E_m)=\lambda(E)$.
\end{proof}
\index{Monge--Amp\`ere measure!multifactor common-measure inequality}

\begin{corollary}\label{cor:bounded-common-measure-mixed}
Let $u,v$ be locally bounded plurisubharmonic functions on an open set $\Omega\subset\mathbb C^n$, $n\geq1$, and let $\mu$ be a positive Radon measure on $\Omega$. If $f,g$ are nonnegative Borel functions on $\Omega$ and
\[
 (\ddc u)^n\geq f\mu,\qquad(\ddc v)^n\geq g\mu,
\]
then, for each integer $1\leq k\leq n-1$,
\begin{equation}\label{eq:bounded-common-measure-mixed}
 (\ddc u)^k\wedge(\ddc v)^{n-k}
                       \geq f^{k/n}g^{(n-k)/n}\mu.
\end{equation}
For $k=0$ and $k=n$ the corresponding inequalities are the hypotheses; for $n=1$ only these occur.
\end{corollary}
\begin{proof}
Apply \cref{thm:bounded-multifactor-common-measure} to $k$ copies of $u$ with density $f$ and $n-k$ copies of $v$ with density $g$.
\end{proof}
\index{Monge--Amp\`ere measure!common-measure mixed inequality}

\begin{corollary}\label{cor:bounded-common-measure-sum}
Under the hypotheses of \cref{cor:bounded-common-measure-mixed},
\[
 (\ddc(u+v))^n\geq(f^{1/n}+g^{1/n})^n\mu.
\]
\end{corollary}
\begin{proof}
By additivity and symmetry of the products, $(\ddc(u+v))^n=\sum_{k=0}^n\binom nk(\ddc u)^k\wedge(\ddc v)^{n-k}$. Bound the terms with $1\leq k\leq n-1$ by \cref{cor:bounded-common-measure-mixed} and the terms with $k=0,n$ by the hypotheses; the binomial formula then gives the assertion. The density on the right is locally $\mu$-integrable, since $f^{k/n}g^{(n-k)/n}\leq f+g$ by the weighted arithmetic--geometric mean inequality, and $f$ and $g$ are locally $\mu$-integrable because $f\mu\leq(\ddc u)^n$ and $g\mu\leq(\ddc v)^n$.
\end{proof}
\index{Monge--Amp\`ere measure!common-measure sum inequality}


\Needspace{10\baselineskip}
\section{Bounded complex Hessians and contact sets}
\label{sec:bounded-complex-hessian-contact}

Equality of two potentials on a measurable set can determine their Hessians there even when that set has no interior. The derivatives used in this assertion are distributional derivatives with measurable densities. We first establish their locality on zero sets.

For an open $D\subset\mathbb R^d$, write $W^{k,2}_{\loc}(D)$ for the locally square-integrable functions whose distributional derivatives of order at most $k$ are locally square-integrable. Mollification converges in this space on smaller open sets, since differentiation commutes with convolution and each of the finitely many derivatives converges in $L^2$.
\index{Sobolev space!local square-integrable derivatives}

\begin{lemma}\label{lem:weak-derivative-zero-set}
If $g\in W^{1,2}_{\loc}(D)$ is real, then $\nabla g=0$ almost everywhere on $\{g=0\}$. Consequently, if $g\in W^{2,2}_{\loc}(D)$, every first and second distributional derivative of $g$ vanishes almost everywhere on that set. Zero sets may be taken for any measurable representative; the assertions are unchanged by null sets.
\end{lemma}
\begin{proof}
For a $C^1$ function $\Phi\colon\mathbb R\to\mathbb R$ with bounded derivative, smoothing $g$ gives the weak chain rule
\[
 \partial_a\Phi(g)=\Phi'(g)\partial_a g.
\]
Indeed, on a relatively compact open set the smoothings converge to $g$ with their first derivatives in $L^2$. The Lipschitz bound gives $L^2$ convergence of their compositions. Pass to a subsequence converging almost everywhere. The difference of the derivative products tends to zero in $L^2$: its term containing the derivative difference is bounded by $\|\Phi'\|_\infty$, while dominated convergence applies to the other term multiplied by the fixed $\partial_a g$. Passing the smooth identity against compactly supported tests proves the formula.

Choose a smooth $\eta$ supported in $(-2,2)$, with $0\le\eta\le1$ and $\eta=1$ on $[-1,1]$, and set $\Phi_\varepsilon(t)=\int_0^t\eta(s/\varepsilon)\,\mathrm ds$. Then $|\Phi_\varepsilon|\le2\varepsilon$, whereas
\[
 \nabla\Phi_\varepsilon(g)=\eta(g/\varepsilon)\nabla g
       \longrightarrow\mathds{1}_{\{g=0\}}\nabla g
       \quad\hbox{in }L^2_{\loc}.
\]
The left side tends to zero distributionally, since $|\Phi_\varepsilon(g)|\le2\varepsilon$. Its $L^2$ limit is therefore zero. This proves the first assertion. For the second, each $\partial_a g$ belongs to $W^{1,2}_{\loc}$ and is zero almost everywhere on $\{g=0\}$. Apply the first assertion to $\partial_a g$ and use the inclusion of these zero sets up to null sets. There are only finitely many derivatives, so their exceptional sets have a null union.
\end{proof}

\begin{lemma}\label{lem:bounded-laplacian-weak-hessian}
Let $f$ be real and continuous on $D\subset\mathbb R^d$. If its scalar distributional Laplacian has a locally bounded measurable density, then $f\in W^{2,2}_{\loc}(D)$.
\end{lemma}
\begin{proof}
Take smoothings $f_\varepsilon$ on a neighborhood of a compact set. They converge locally uniformly to $f$, and $\Delta f_\varepsilon\to\Delta f$ in $L^2$ there. If $\chi$ is a smooth compactly supported real cutoff, integration by parts gives, for a real smooth $h$,
\[
 \begin{aligned}
 \int\chi^2|\nabla h|^2
 &=-\int\chi^2h\Delta h-2\int\chi h\nabla\chi\cdot\nabla h\\
 &\le \tfrac12\int\chi^2|\nabla h|^2
       +\int\chi^2|h\Delta h|+2\int h^2|\nabla\chi|^2.
 \end{aligned}
\]
Apply this to $h=f_\varepsilon-f_\delta$. Uniform convergence and the local $L^2$ bound for its Laplacian show that the gradients form a Cauchy family in $L^2$ on every smaller compact set. Their limits are the distributional first derivatives of $f$.

Now choose $\chi=1$ near the compact set under consideration. The functions $v_\varepsilon=\chi f_\varepsilon$, extended by zero, are smooth and compactly supported, and
\[
 \Delta v_\varepsilon
 =\chi\Delta f_\varepsilon+2\nabla\chi\cdot\nabla f_\varepsilon
                                      +f_\varepsilon\Delta\chi
\]
is Cauchy in $L^2(\mathbb R^d)$. For every real smooth compactly supported $v$, twice integrating by parts gives
\[
 \sum_{a,b=1}^d\|\partial_a\partial_b v\|_2^2
 =\sum_{a,b=1}^d\int(\partial_a^2v)(\partial_b^2v)
 =\|\Delta v\|_2^2.
\]
Apply this identity to $v_\varepsilon-v_\delta$. All real second derivatives are Cauchy in $L^2$; passing against tests identifies their limits with the distributional derivatives of $\chi f$. Where $\chi=1$ these are the derivatives of $f$. A compact exhaustion proves the assertion.
\end{proof}

Let $\Omega\subset\mathbb C^n$. In the next results a continuous function has a bounded complex Hessian if its distributional coefficients $f_{j\bar k}=\partial^2 f/\partial z_j\partial\bar z_k$ have locally bounded measurable densities. Such a function satisfies the preceding lemma because the scalar Laplacian is $4\sum_jf_{j\bar j}$.
\index{Hessian!bounded complex coefficients}

\begin{proposition}\label{prop:bounded-complex-hessian-contact}
If real continuous $f,g$ belong to $W^{2,2}_{\loc}(\Omega)$, then their real weak first and second derivatives agree almost everywhere on $E=\{f=g\}$. In particular, their weak complex-Hessian coefficients satisfy
\[
 \mathds{1}_E f_{j\bar k}=\mathds{1}_E g_{j\bar k}
 \quad\hbox{almost everywhere for all }j,k.
\]
\end{proposition}
\begin{proof}
Apply \cref{lem:weak-derivative-zero-set} to $f-g$. Complex derivatives are fixed linear combinations of the real derivatives, giving the formula. In particular the result applies to continuous functions with bounded scalar Laplacians, by \cref{lem:bounded-laplacian-weak-hessian}, and to functions with bounded complex Hessians.
\end{proof}

For Hermitian $n\times n$ matrices $A_1,\dots,A_n$, let $D(A_1,\dots,A_n)$ be $1/n!$ times the coefficient of $t_1\cdots t_n$ in $\det(\sum_jt_jA_j)$. Thus $D(A,\dots,A)=\det A$; in dimension zero put $D()=1$.

\begin{proposition}\label{prop:bounded-complex-hessian-density}
Let $u_1,\dots,u_n$ be continuous plurisubharmonic functions with bounded complex Hessians on $\Omega$. With $A_j=((u_j)_{a\bar b})$, their Bedford--Taylor measure is
\[
 \ddc u_1\wedge\cdots\wedge\ddc u_n
       =\frac{n!}{\pi^n}D(A_1,\dots,A_n)\,\mathrm dV_{2n}.
\]
If real continuous $v_1,\dots,v_n$ belong to $W^{2,2}_{\loc}(\Omega)$, put $B_j=((v_j)_{a\bar b})$ for their weak complex Hessians and $E=\bigcap_j\{u_j=v_j\}$. Then
\[
 \mathds{1}_E\ddc u_1\wedge\cdots\wedge\ddc u_n
       =\frac{n!}{\pi^n}\mathds{1}_E D(B_1,\dots,B_n)\,\mathrm dV_{2n}.
\]
The restricted density on the right is locally bounded even if the $v_j$ are not plurisubharmonic; no integrability of $D(B_1,\dots,B_n)$ off $E$ is needed. The empty product in dimension zero is counting measure, and both identities have that interpretation.
\end{proposition}
\begin{proof}
For smooth $u_j$, polarize \eqref{eq:bt-determinant-normalization}. The coefficient of $t_1\cdots t_n$ on the exterior-product side is $n!\ddc u_1\wedge\cdots\wedge\ddc u_n$, giving the asserted normalization. For continuous $u_j$, take their local radial smoothings. They converge uniformly on compact subsets, so \cref{prop:bt-uniform-continuity} gives convergence of their products. Their complex-Hessian coefficients are the convolutions of the bounded densities $A_j$. At the common Lebesgue points they converge to $A_j$, and they are uniformly bounded near each test support. The polynomial $D$ therefore converges in $L^1$ there by dominated convergence. This identifies the limiting measure and proves the first formula. By \cref{prop:bounded-complex-hessian-contact}, $A_j=B_j$ almost everywhere on $\{u_j=v_j\}$, hence on $E$. Substitution proves the second formula. In dimension zero there are no Hessian coefficients or exceptional derivative sets, and the convention gives the assertions directly.
\end{proof}

Smooth local potentials of closed real background forms may be added to both members of each contact pair. The last identity consequently applies in a holomorphic chart when the plurisubharmonic potentials have bounded complex Hessians and their contact barriers have real weak $W^{2,2}$ derivatives. Its weak-derivative argument does not require an almost-everywhere classical real Hessian.


\Needspace{30\baselineskip}
\section{Gradients against positive currents}
\label{sec:bt-current-gradients}
The following local chain rule also applies to the cutoff arguments on compact manifolds. Here $\mathrm d^{\mathrm c}=(\partial-\bar\partial)/(4\pi\mathrm i)$, so $\mathrm d\mathrm d^{\mathrm c}=\ddc$.

\Needspace{23\baselineskip}
\begin{lemma}\label{lem:bounded-psh-current-chain-rule}
Let $\Theta$ be a closed positive current of bidimension $(1,1)$ on a coordinate domain, and let $v$ be locally bounded psh. The current $\mathrm d(v\Theta)$ has measure coefficients, denoted $\mathrm dv\wedge\Theta$. There is a positive gradient measure
\begin{equation}\label{eq:bounded-current-gradient-measure}
 E_\Theta(v)=\mathrm dv\wedge\mathrm d^{\mathrm c}v\wedge\Theta
       =\tfrac12\ddc(v^2\Theta)-v\,\ddc v\wedge\Theta.
\end{equation}
Smooth decreasing approximations have gradients converging strongly in the coefficient-matrix energy seminorm on compact subsets. For every smooth real $f$, the bounded products satisfy
\begin{equation}\label{eq:bounded-current-chain-rule}
 \mathrm d(f(v)\Theta)=f'(v)\,\mathrm dv\wedge\Theta.
\end{equation}
For a compactly supported smooth real one-form $\eta$, a cutoff $\tau\geq0$ equal to one near its support, and a bounded Borel function $a$, there is a constant $C_\eta$ such that
\begin{equation}\label{eq:current-gradient-cauchy-schwarz}
 \left|\int a\eta\wedge\mathrm dv\wedge\Theta\right|^2
 \leq C_\eta\left(\int\tau E_\Theta(v)\right)
                 \left(\int_{\operatorname{supp}\eta}|a|^2\,
                                         \Theta\wedge\beta\right).
\end{equation}
\end{lemma}
\begin{proof}
All assertions are local near a compact test support. Add a constant to make $v$ and its smooth radial approximations $v_l\downarrow v$ nonnegative and uniformly bounded there. For a smooth real function $w$, the differential product rule gives
\[
 \mathrm dw\wedge\mathrm d^{\mathrm c}w\wedge\Theta
       =\tfrac12\ddc(w^2\Theta)-w\ddc w\wedge\Theta.
\]
The measure on the left is positive. Applied to $w=v_l-v_m$ and tested with $\tau$, this identity gives energy tending to zero as $l,m\to\infty$. Indeed \cref{cor:bt-polynomial-weights} makes the squared-weight term tend to zero after integration by parts against $\ddc\tau$, and each of the four expanded Hessian terms tends to the same integral with weight $v$ and Hessian $\ddc v$. Their signs cancel. For two independent indices the convergence is uniform on tails: otherwise select violating pairs with both indices strictly increasing and apply that same decreasing-continuity theorem.

Here is the meaning of the energy seminorm, including when $\Theta$ is singular. Its positive coefficient matrix has bounded Radon--Nikodym entries relative to $\mu=\Theta\wedge\beta$, by \cref{thm:positive-current-order-zero}. Write that measurable Hermitian matrix as $H$. For this bidimension, its quadratic form is the density of $\mathrm i\xi\wedge\bar\xi\wedge\Theta$ for a constant $(1,0)$-form $\xi$. Positivity for a countable dense family of $\xi$ makes $H$ positive semidefinite almost everywhere. In coordinates the gradient energy is the squared $L^2(\mu)$ norm of $H^{1/2}\partial w$, with the fixed normalization absorbed into $H$. The preceding estimate makes these vectors for $v_l$ Cauchy on compact subsets. Their strong limit defines the gradient modulo the zero seminorm. Multiplication by a bounded measurable scalar respects that quotient. Pointwise matrix Cauchy--Schwarz bounds the derivative paired with $a\eta$ by the two norms in \eqref{eq:current-gradient-cauchy-schwarz}. It also bounds every derivative coefficient in variation. Distributional convergence of $v_l\Theta$ to $v\Theta$, by dominated convergence, identifies that measure derivative with $\mathrm d(v\Theta)$.

The squared norms converge in $L^1(\mu)$, so their limit is a positive measure. Polynomial weighted continuity identifies its distribution with the right side of \eqref{eq:bounded-current-gradient-measure}. This also proves independence of the smooth approximations. Finally $f(v_l)\Theta\to f(v)\Theta$ by dominated convergence. The products $f'(v_l)$ times the gradients converge strongly: subtract the limiting gradient first, then use dominated convergence on its squared energy with $f'(v_l)\to f'(v)$. Passing through the smooth chain rule proves \eqref{eq:bounded-current-chain-rule}. A constant shift only translates the argument of $f$ and does not change the gradients, completing the proof.
\end{proof}

\index{gradient measure!positive currents}\index{chain rule!bounded psh potentials}

